\documentclass[11pt]{article}
\usepackage[a4paper,margin=28mm]{geometry}
\usepackage{amsmath,amssymb,amsthm,mathtools}
\usepackage[T1]{fontenc}
\usepackage{lmodern}
\usepackage{microtype}
\usepackage{xurl}
\usepackage[hidelinks]{hyperref}
\newtheorem{theorem}{Theorem}
\newtheorem{lemma}[theorem]{Lemma}
\newtheorem{proposition}[theorem]{Proposition}
\newtheorem{corollary}[theorem]{Corollary}
\newtheorem{definition}[theorem]{Definition}
\newcommand{\F}{\mathbb F}
\newcommand{\cost}{\operatorname{cost}}
\newcommand{\id}{\operatorname{id}}
\newcommand{\polylog}{\operatorname{polylog}}
\title{Ordinary Binary Polynomial Multiplication below $n\log n$\\
\large Controlled tensor routing on fixed tapes}
\author{\'Arp\'ad Csizmadia\\
\small\href{mailto:arpad.csizmadia.public@gmail.com}{\texttt{arpad.csizmadia.public@gmail.com}}}
\date{31 October 2026\\ \small Version 0.1.0-preprint. Apache-2.0.}
\hypersetup{pdfauthor={\'Arp\'ad Csizmadia},pdftitle={Ordinary Binary Polynomial Multiplication below n log n}}
\begin{document}
\maketitle
\begin{abstract}
We give a transfer from a fast equal-chunk tape interchange
to ordinary dense multiplication in $\F_2[X]$. The main new interfaces
are tensors of permutations on unequal even radices and a coupled
cyclic-index to prime-residue index permutation. A flat perfect shuffle
differs from its mixed-radix tensor by a low-endpoint cyclic rotation.
A ragged parity-splitting network keeps the original even dimensions
throughout its recursion. Controlled versions use completed target-only
kernels and a temporary control gauge. A balanced CRT tree then needs
only logarithmically many controlled tensor stages. The cyclotomic
coefficient word can be enlarged to pay for polynomial-time address
metadata without changing its small-word arithmetic regime. Together
with the characteristic-two local Fourier compiler and the explicit
interchange and endpoint-kernel constructions in the appendices, this gives
the bound $O(n(\log n)^{1-2^{-38}})$ in the deterministic
fixed-finite-tape bit model. This is a research preprint open for mathematical review;
no external verification or priority is claimed.
\end{abstract}

\paragraph{AI assistance and formal verification.}
This work was developed with substantial generative AI assistance in
mathematical exploration, algorithm and proof development, writing,
coding, and literature review. OpenAI Codex~\cite{CodexTool26}
assisted release preparation and source inspection. Historical
model versions were not consistently recorded; no complete
model-level attribution is claimed. Complete Lean formalization is unfinished
and planned for a later update. This manuscript is a research
preprint open for mathematical review.

\paragraph{Status and scope.}
The input and output are ordinary monomial coefficient lists over
$\F_2$. The function is ordinary polynomial multiplication. We count
all address processing, padding, scalar preparation and record movement.
The finite motif, interchange, endpoint-kernel and ordinary-I/O proofs are
expanded in this manuscript and the selected proof-support records.
The whole composition must be read as an actual tape construction.
Finite checks below corroborate identities; they do not prove asymptotic
running times. The companion manuscript, \emph{Ordinary Polynomial
Multiplication over Every Fixed Finite Field}, is provided alongside
this paper in the review release.

\paragraph{Relationship to the companion papers.}
The accessible companion, \emph{Exact Binary Polynomial Multiplication
from Fast Pascal Transforms}, develops binary prefix/Taylor evaluation
and a trace-pairing residue decoder. The present construction shares
finite-network and equal-chunk interchange antecedents but uses a
different outer route: a cyclotomic host, coefficient CRT components,
a balanced CRT routing tree and ordinary cyclic-product reconstruction.
It does not invoke the Pascal companion's complete multiplier as a
subroutine. The general finite-field companion combines the binary
bound proved here with its odd-characteristic construction. The present
proof does not depend on that combined consequence, so this relationship
is not circular.

\paragraph{Prior work.}
An earlier fixed-prime benchmark is
\[
 O_p(n\log n\,8^{\log^* n}),
\]
due to Harvey, van der Hoeven and Lecerf~\cite{HVHL17}.
Harvey and van der Hoeven supply the unconditional cyclotomic host and
an $O(n\log n\,4^{\log^*n})$ fixed-prime polynomial bound
\cite{cyclomult}. Their $O(n\log n)$ finite-field theorem treats the uniform-field
problem under an explicit near-one Linnik hypothesis \cite{ffnlogn}.
The present fixed-binary-field construction builds on their
unconditional cyclotomic host and supplies ordinary coefficient
reconstruction within its stated tape cost.
Harvey's October 2026 discussion identifies fixed-field polynomial
multiplication as a separate problem from the new integer claims
\cite[Section 4]{Harvey26Context}.
Van der Hoeven and Larrieu exploit Frobenius symmetry in finite-field
Fourier transforms, with applications to binary polynomial multiplication
\cite{Frobenius17}. Their extension arithmetic and conversion costs
are part of that analysis. Our stated all-length bound requires the
complete ordinary-input/output and routing accounting developed here.
The 2026 interchange preprint supplies the finite-network tape method
\cite{swap}; its global integer multiplication theorem is not needed.
The exact-complex Fourier paper supplies the local Newton/Toeplitz and
borrowed-coordinate mechanism \cite{local}; its complex scalars and
unit-cost indexing do not give a finite-field bit theorem directly.
It also states an explicit exact complex-polynomial multiplication
corollary in that arithmetic model. The corollary retains its
specified-root and address-cost assumptions and expressly excludes
finite fields and bit complexity. Our claimed coefficient field and
machine model therefore require the separate transfer proved here.
Cantor--Kaltofen \cite{cantor-kaltofen} give classical formal
convolution circuits over arbitrary coefficient algebras. We print
an elementary characteristic-two triadic circuit with looser
secondary logarithms to make our required generation bound explicit.
Arbitrary-size permutation networks have substantial prior art,
including \cite{benes}. We make no novelty claim for their abstract
existence. Our required statement is a priced tensor execution on
the existing mixed-radix tape, including controls and dummies.
Current follow-up work also studies larger integer-multiplication
exponents \cite{integer-followups}; a separate repository states a
finite-machine theorem for multiplication of natural numbers
\cite{integer-lean}. We do not use those claimed end-to-end bounds
as premises. The cited primary and related sources are listed in the
bibliography; this is not an exhaustive priority determination.

Together these works provide the host arithmetic, finite-network
methods and transform mechanisms used in this paper. The specific
advance claimed here is an ordinary binary polynomial multiplier
with a fixed positive saving below the logarithmic exponent one,
obtained by completing the controlled mixed-radix and CRT routing
transfer. The preceding fixed-field bound and the conditional
uniform-field theorem explain why that threshold is a meaningful
target. The tiny saving is a theoretical asymptotic statement;
the earlier algorithms and their structural contributions remain
the foundations of this construction.

\section{Model and supplied kernels}

One machine has a fixed finite alphabet and a fixed finite number of
one-dimensional tapes. Heads move by one cell per step. Shape parameters
are ordinary binary integers. Arrays occur in lexicographic order.
A record has $D$ arbitrary bits, and $V$ denotes the entire bit volume.
Public addresses are positions of records, never payload controls.

We reconstruct the following interchange and endpoint kernels in the
appendices, with further implementation detail in the companion dossier.
Their dependency direction is explicit: the rational-frame interchange
uses the binary scalar motif and ordinary scans; the endpoint kernel
then uses interchange and the finite ternary-phase word. The
interchange proof never invokes H2 or a polynomial multiplier.

\begin{proposition}[Interchange and native endpoint kernels]
\label{prop:suppliers}
There is a fixed $\tau=1-2^{-36}<1$ such that equal address chunks
in $[P]\times[2^u]\times[G]\times[2^u]\times[B]$ can be interchanged in
$O(Vu^\tau)$ bit time, uniformly in positive $P,G,B$.
Put $\gamma=1/(2-\tau)<1$; aggregate rectangular exchanges have cost
$O(VL^\gamma)$ for total address width at most $L$.

Let $n_i$ be even native radices, $M=\prod_i n_i$,
$L=\lceil\log_2(2M)\rceil$, and
\[
 L^3\le n_i\le R,\qquad \log R=O((\log L)^c).
\]
On each adjacent pair $(2a,2a+1)$ put
$H=\left(\begin{smallmatrix}1&1\\1&0\end{smallmatrix}\right)$.
The tensor of these canonical native $H$ maps has cost
$O(MD L^\gamma\operatorname{polylog}L)$, uniformly for external $D\ge1$.
All private padding and work-tape returns are included.
\end{proposition}

Appendix~\ref{app:swap} supplies interchange directly, re-instantiating
the method of \cite{swap}; Appendices~\ref{app:finite-phase} and
\ref{app:native-endpoint} prove the endpoint consequence. Thus the fast
interchange statement is reconstructed here, rather than left as an
unproved existence premise of the multiplication transfer.

The native kernel uses the quotient labels as carry guards, rather
than individually padding the native radices. A frozen endpoint mask
on $u$ digits is a constant word of weighted integer rotations plus
exact sparse boundary repair. A large untouched suffix supplies a dock
and exponentially many opaque lanes. Its numerical padding survives
the complete binary kernel, including every internal $\F_{16}$ letter.
The supplier proof must retain those lifetimes.

\paragraph{Deriving the aggregate rectangle bound.}
For completeness, interchange two adjacent binary fields of widths
$a,b$, with $a+b\le L$. Choose an integer threshold
$u=\lceil L^{1/(2-\tau)}\rceil$. While both residual fields have
width at least $u$, use the block-exchange Euclidean algorithm.
If $a\le b$, split the second field into its first $a$ bits and the
remaining $b-a$, exchange the two equal $a$-bit fields, and retain
the already placed prefix. The other case is symmetric. Each step
removes its shorter width $v$ from the unresolved total. Thus the
sum of those $v$ is at most $L$, and, since $\tau-1<0$,
\[
 \sum_{v\ge u}O(Vv^\tau)
   \le O(Vu^{\tau-1})\sum v
   =O(VLu^{\tau-1}).
\]
When one residual field has width below $u$, move its bits individually
through the other field, at $O(V)$ per bit, using streamed two-row
chunk interleaves. Each suffix block is copied as a block, with
amortized counters. The last cost is $O(Vu)$.
Move a right-hand short field's bits first to last to the front of
the long field; move a left-hand short field's bits last to first
after it, retaining already placed bits. Both routes preserve the
short field's internal order.
Both terms are $O(VL^\gamma)$ for $\gamma=1/(2-\tau)$.
All calls keep the same full volume $V$ and already placed fields are
spectators. For rational fixed $\tau$, an integer power comparison
computes the threshold; its polynomial-in-$L$ setup fits in the
array volume. To exchange native aggregate ranges, pad each whole
range once to a binary range (combined factor below four), perform
this exchange and crop the nested valid rectangle after completion.
There is no separate pad for each native axis.

\section{A tensor perfect shuffle}

Write $n_i=2k_i$ with the first digit most significant. Define
\[
 S_i(k_i h_i+\ell_i)=2\ell_i+h_i,
 \quad h_i\in\{0,1\},\quad 0\le\ell_i<k_i.
\]
Also define $d_M(X)=2X\bmod(M-1)$ for $X<M-1$ and fix $M-1$.
Splitting the flat stream into its two halves and interleaving records
implements $d_M$ in $O(V)$ time.

\begin{lemma}[Post-shuffle endpoint correction]
\label{lem:shuffle}
The tensor $\bigotimes_i S_i$ and its inverse have cost
$O(VL^\gamma\polylog L)$ under Proposition~\ref{prop:suppliers}.
\end{lemma}
\begin{proof}
For $x_i=k_i h_i+\ell_i$ and incoming carry $c\in\{0,1\}$,
\[
 2x_i+c=n_i h_i+(2\ell_i+c).
\]
The outgoing carry is exactly $h_i$, independent of $c$.
End-around doubling therefore has output
$y_i=2\ell_i+h_{i+1}$, cyclically in $i$.
Consequently
\[
 \bigotimes_i S_i=r\circ d_M,\qquad
 r(y)_i=2\lfloor y_i/2\rfloor+(y_{i-1}\bmod2).
\]
This rotates low endpoint bits \emph{after} the flat shuffle; it does
not require a cheap rotation of the original high-half bits.

Index the axes from zero. Reverse endpoint bits on $[0,t-1)$ and then
on $[0,t)$. A bit at $i<t-1$ moves first to $t-2-i$ and then to $i+1$;
the last bit moves to zero. This is precisely $r$. For a reversal
$[a,b)$, a left-half interval $I=[c,d)$ pairs with
$J=[a+b-d,a+b-c)$, in reverse local order. Partition the left half
into such intervals of length $u\le h$, where
\[
 \ell_R=\lceil\log_2R\rceil,\qquad
 h=\min\{\lfloor t/16\rfloor,\lfloor L/(40\ell_R)\rfloor\}.
\]
The complement of two selected intervals has at most three linear
intervals. Since $t\ge16h$, one has at least $u$ axes; choose a
contiguous $u$-axis tile auxiliary $W$ there.

Place the two selected banks $X,Y$ and $W$ at the front in that order.
Extract them successively from the remaining axis list: at each step
exchange the whole preceding aggregate gap with the requested interval.
At most three adjacent aggregate rectangle exchanges suffice. All
unselected axes keep their relative order; reverse the same exchanges
after the batch. Each exchange crops its two aggregate pads before the
next one and costs $O(VL^\gamma)$. No endpoint gather is performed.

With $\pi(j)=u-1-j$, the chronological endpoint masks
$Y_{\pi(j)}\mathrel{\oplus}=X_j$,
$X_j\mathrel{\oplus}=Y_{\pi(j)}$,
$Y_{\pi(j)}\mathrel{\oplus}=X_j$
exchange the paired bits. Each reads only the opposite frozen bank
and uses $W$ as its separate dirty tile auxiliary; its weighted offset
uses the target bank's own native place values. Every quotient label
is fixed. The three banks use at most $3L/40$ address bits. After
padding the one remaining suffix to $2^s$ and borrowing a dock of
width $c\le L/40+1$, its low index supplies
$R_{\rm op}=2^{s-c}D\ge2^{9L/10-O(1)}D$ opaque bits per record.
This pays the weighted-mask metadata and exact boundary repair even
when $D=1$. There are no phase-role bits in these direct masks.
The completed masks restore the dock and fix the low suffix index,
so the suffix pad is cropped before the inverse layout.

Each mask costs $O(V(u\log R)^\tau+V)$.
Since $\gamma>\tau$ and $h\log R\le L/40$, a batch costs
$O(VL^\gamma)$. From $L-2<\log_2M\le t\log_2R$,
$h\ge L/(80\ell_R)$ for sufficiently large admissible $L$.
Together with $t\le L/(3\log_2L)$ this gives $O(\log R)$ batches.
Summing proves the bound. Reversing the actual word gives the inverse;
the bounded remaining range uses a direct finite branch.
\end{proof}

\section{Permutations on exact even dimensions}

An adjacent pair can select either $I$ or the exchange matrix $J$.
Over $\F_4$, both have a common word of exactly eight $H$ letters
interleaved with unit diagonals. Tensoring a slot uses the supplied
canonical kernel. The diagonal coefficient factors by local axes;
its lexicographic list is generated by prefix products in $O(M)$
fixed-field operations. The completed word is binary. Its lane pad
can be removed only at that completed boundary.

Odd-start adjacent matchings are conjugates by coordinatewise $+1$.
Explicitly, if $T(x)=x+1\bmod n$, an odd-start pair $(i,i+1)$ is
$TKT^{-1}$ with $K$ switching $(i-1,i)$. Set $K=I$ on every
irrelevant pair, in particular $(n-2,n-1)$ whose conjugate would
wrap to $(n-1,0)$. Thus the completed matching has no cross-boundary
or wrap switch, even when its temporary $K$ crosses a ragged parent
boundary.
At a synchronized odd-start stage, translate every axis in the fixed
target set, including axes whose matching is identity. Choose $K=I$
on those axes; conjugation then cancels their translations. Thus the
translated set, flat offset and source mask are independent of switch
settings, including any control address in the controlled compiler.
This selective translation costs $O(V)$: flat rotation by the sum of
the selected native place values agrees except when a selected digit
wraps. The exceptional fraction is at most $t/L^3\le L^{-2}$.
Rotate the good-source mask with the payload. For bad output slots,
compute $x=P^{-1}(y)$ and the correct destination $T(x)$, sort those
records by destination, and refill the holes. Since $P=T$ on good
sources and both are permutations, $P(B)=T(B)$. This does not assume
that the bad source set is invariant. Keys, record sorting and
schoolbook decoding cost
\[
 O(\delta ML(D+L)+\delta ML^2)=O(V).
\]
Here the $L^2$ exception-address bound does not perform $t$ full
$L^2$ divisions. If $b_i=\lceil\log_2n_i\rceil$, then
$\sum_i b_i\le2L$. Streaming a current at-most-$L$-bit numerator
through a $b_i$-bit remainder register divides by $n_i$ in
$O(Lb_i+L+b_i)$ tape steps. Summing gives $O(L^2)$ for all digits;
the weighted Horner reconstruction has the same bound.

The source-good mask also has a linear implementation. Starting with
one bit, append native factors in their physical order. For each bit
$a$ of the previous mask, print $a$ for the $n_i-1$ nonmaximal values
and zero at the maximal value if axis $i$ is selected; otherwise
print $n_i$ copies of $a$. Counter reset work is at most its ensuing
$n_i$-bit output block. Since the intermediate mask lengths grow by
factors at least two, their sum is below $2M$. Rotate that mask with
the flat data permutation, so it marks exactly the output holes $P(B)$.
Compute and copy a full key only at a bad record. A monotone output
counter otherwise advances amortized in $O(M)$ total time. Stable
binary radix sorting in $L$ passes costs
$O(KL(D+L))$ for $K\le\delta M$ bad records, including the full
payload and header scans. In each pass stage one fixed-width item,
remember its selected key bit, and copy the item to the corresponding
stable zero/one bucket. Concatenate the buckets and rewind; a constant
number of tapes and $O(K(D+L))$ steps suffice per pass, including
staging and returns. Binary counters have amortized linear cost over
each scan; their per-item resets cost at most that item's width.
Since the sorted destinations are exactly
$P(B)=T(B)$, refill bad holes in that order; there is no $L$-bit key
comparison at each good position. This establishes the stated bound
even for $D=1$. The negative translation uses nonzero-digit source
masks and the same proof.

\begin{lemma}[Ragged global-shuffle network]
\label{lem:ragged}
Every permutation of an original even $n$ has a word of $O(\log n)$
levels consisting of adjacent switches and the \emph{same} full $S_n$
and $S_n^{-1}$. No recursive odd-sized child is padded.
\end{lemma}
\begin{proof}
At a level, the residual permutation preserves \emph{each} interval
setwise in a public ordered partition into contiguous parents.
Pair positions from each parent's
start, leaving its last position unmatched if its size is odd.
Use the same output pairing. The bipartite multigraph has input
pair/singleton vertices, output pair/singleton vertices, and one edge
per packet. Even parents give even cycles. An odd parent gives one
path between its two singleton vertices, plus even cycles. That path
has an odd number of edges. Its endpoint edge colors can both be
forced to the global parity of the parent's last position.

Alternating edge colors give input switches $E_{\rm in}$ and output
switches $E_{\rm out}$ so that the residual has identity
\[
 \pi=E_{\rm out}S_n C S_n^{-1}E_{\rm in}.
\]
The full inverse shuffle stably splits global even and odd positions.
For original even $n$, write this split as
\[
 p_n(x)=S_n^{-1}(x)=
 \begin{cases}
 x/2,&x\text{ even},\\
 n/2+(x-1)/2,&x\text{ odd}.
 \end{cases}
\]
A parent interval $[a,b)$ therefore has the two image fragments
\[
 \left[\left\lceil a/2\right\rceil,\left\lceil b/2\right\rceil\right),
 \qquad
 \left[n/2+\left\lfloor a/2\right\rfloor,\,
       n/2+\left\lfloor b/2\right\rfloor\right).
\]
The even fragments concatenate to $[0,n/2)$ and the odd fragments
to $[n/2,n)$, skipping empty fragments.
The new contiguous block list is
\[
 [\text{all parents' even fragments}]
 [\text{all parents' odd fragments}].
\]
The endpoints depend only on $n$ and the parent partition, independently
of $\pi$ and its switch colors. The residual $C$ preserves each new
fragment setwise. The recursive call is on this single $n$-wire
residual with the whole updated partition. Thus each level adds only
a fixed number of slot types to one recursive word.
Track this order; do not gather siblings. Each fragment has size at most
the ceiling of half its parent. After $\lceil\log_2n\rceil$ levels all
blocks are singletons. Each adjacent switch stage splits into its
even-start and odd-start families. Short axes can continue passive
singleton levels; their remaining shuffles cancel on unwind.
\end{proof}

Synchronize the networks for all axes to $\lceil\log_2R\rceil$ levels.
Lemmas~\ref{lem:shuffle}--\ref{lem:ragged} and the adjacent-switch
words give a generic local-permutation tensor cost
\begin{equation}
 \label{eq:generic}
 O(VL^\gamma\polylog L).
\end{equation}
This retains the parameter of the supplied kernels: each adjacent-switch
word uses a fixed number of native-$H$ calls, each shuffle has the
same $\gamma$ tariff, and there are $O(\log R)$ levels. Translations
and fixed-field diagonal scans cost $O(V)$. Thus synchronization adds
only a fixed power of $\log L$. Public local tables and network
settings have polynomial-time sequential generation in $R,L$.

\section{Controlled tensors and the control gauge}

Partition axes into targets $T$ and controls $C$. For each control
address $c$ let $\pi_i(c)$ be a permutation on target axis $i$.
Controls are address coordinates, not input-dependent payload values.

\begin{lemma}[Controlled compiler]
\label{lem:controlled}
Let $G(R,L)\ge1$ bound the time to generate one outer $\F_4$ scalar
coefficient at a current record address, including address decoding,
control-gauge inversion, all required local tables and ragged settings.
Assume this is a fixed polynomial-time sequential procedure with fixed
tapes. Then the controlled map
\[
 (c,x_T)\longmapsto(c,(\pi_i(c)(x_i))_{i\in T})
\]
has bit cost
\[
 O\bigl((MDL^\gamma+MG(R,L))\polylog L\bigr)
\]
with fixed tapes and ordinary order. In particular $G(R,L)=O(D)$
gives \eqref{eq:generic}. For the explicit CRT families below,
$G(R,L)=O((R+L+2)^{12})$; thus $D\ge(R+L+2)^{64}$ suffices.
\end{lemma}
\begin{proof}
First construct the completed subset kernel
$K_T=H_T\otimes I_C$. Use a static common eight-$H$ word for $H$ versus
$I$, tensoring each inner letter across all axes. The control axes may
move \emph{inside} this implementation. At its completed boundary,
$K_T$ fixes them. Insert controlled scalar stages only between completed
$K_T$ calls. The outer eight-$H$ switch word is then valid on each
fixed-control fiber. Its coefficient is
$\kappa(x,c)=\prod_{i\in T}\kappa_i(x_i,c)\in\F_4^*$.
We do not assume factorization by control axes.

The fragment-endpoint formulas in Lemma~\ref{lem:ragged} make the
partition and shuffle topology independent of the control address.
Only the switch settings depend on $c$, so a single fixed slot
schedule applies to all fibers.
Use the full tensor shuffle on \emph{all} axes. Give $C$ identity
switches. At an outer slot their temporary physical value is
$q_e(c)=S_C^e(c)$, where $e$ is the signed number of completed forward
minus inverse shuffles. At each scalar stage enumerate the current
native lexicographic addresses and decode logical
$c=q_e^{-1}(x_C)$. Generate the target settings for this logical $c$.
At every odd-start stage use the same translation of all axes in the
fixed set $T$, with identity $K_i(c)$ on idle axes or fibers. The flat
approximation can temporarily carry into a control coordinate on bad
records. Complete its exact repair before reading any control-dependent
setting. Completed subset kernels and translations then preserve the
physical controls; a shuffle increments $e$ in its prescribed direction.
Induction over the word proves the fiber invariant. Descent and unwind
balance, so $e=0$ at exit and the controls are restored.

For the CRT affine families, one concrete generator decodes addresses
by schoolbook division, computes the control CRT sums, generates the
at-most-$R$ local table, and constructs the pair graph by sequential
scans. All axes, levels, records, and full-list lookups are bounded in
the companion ledger by $O((R+L+2)^{12})$ time. Fix the dominating
polynomial $P_1=(R+L+2)^{64}$ before choosing $D$. For each record park
its $D$ bits locally, generate one coefficient, apply its fixed two-bit
map, and advance. In general this scalar stage costs $O(M(D+G))$;
visited work intervals and returns are included in the generator bound.
There are only $\polylog L$ outer stages, proving the displayed
general bound. The CRT generator just described has $G=O((R+L+2)^{12})$,
which the chosen word width dominates, so its scalar stage costs $O(V)$.
No whole-array scan is nested inside a per-record lookup. This also
aligns coefficients to the actual post-shuffle stream order.
\end{proof}

\section{Enlarging the cyclotomic word}

\begin{lemma}[Uniform host initialization]
\label{lem:host-initialization}
From the binary input length $n$, for all sufficiently large $n$ a
deterministic fixed-tape procedure constructs a $2$-admissible tuple
$(q_0,\ldots,q_e)$, its product $2n<N<4n$, the attached common-multiple
data, and $\alpha,\Phi_\alpha,\rho$ with $q_0\mid\alpha$ and
$\rho$ a principal $(N/q_0)$-root. Its bit time is $N^{o(1)}$,
including writing and verifying the factor and root data.
\end{lemma}
\begin{proof}
Apply \cite[Proposition 3.4]{cyclomult} at $p=2$ and target $2n$;
its tuple and auxiliary data are part of the constructed admissible
length. Apply Proposition 3.8 of that source to this tuple. Their
respective time bounds are
$2^{O((\log\log n)^2)}$ and $2^{O((\log\log N)^4)}$ at fixed $p=2$.
Both are $N^{o(1)}$ in the source's deterministic fixed-tape model.
Write $L_0=\lceil\log_2N\rceil$ and
$\ell=\lceil\log_2L_0\rceil$; its divisor has
$\log\alpha=O(\ell^4)$.

Verification uses only polynomial-size setup operations. Trial division
checks the emitted prime list and its product, since each prime is at
most $R_0=2^{\ell^2}$; monic polynomial arithmetic checks the generated
cyclotomic data. Put $K=N/q_0$. Check $\rho^K=1$ and
\[
 \gcd(\rho^{K/q_i}-1,\Phi_\alpha)=1\qquad(1\le i\le e).
\]
The squarefree quotient is a product of fields. These tests force
order $K$ in every component, hence principality by the geometric-sum
identity. Modular powers, gcds and all certificates cost a polynomial
in $\alpha,R_0,L_0$, which is $2^{O(\ell^4)}=N^{o(1)}$.
Schoolbook arithmetic suffices in this setup; it uses no call to the
proposed multiplier.
\end{proof}

If $n$ is implicit in end-marked coefficient lists, counting and copying
them costs $O(n)$ by amortized counters. The bounded initial length
range uses the direct branch; no factorization is supplied as advice.

Use the unconditional host of \cite{cyclomult}. Put
$L_0=\lceil\log_2N\rceil$, $\ell=\lceil\log_2L_0\rceil$ and
$R_0=2^{\ell^2}$. Its squarefree $N$ has prime factors in
$(L_0^3,R_0]$. The constructed divisor $\alpha$ uses $O(\ell^2)$
factors and satisfies
\[
 \log\alpha=O(\ell^4),\qquad
 -\log\frac{\varphi(\alpha)}\alpha=O(\ell^2/L_0^3).
\]
The distinguished prime $q_0$ divides $\alpha$. Proposition 3.8
constructs $\rho$, a principal $(N/q_0)$-root in
$\F_2[Y]/\Phi_\alpha(Y)$, as well as this cyclotomic polynomial.
Then $\omega=\rho^{\alpha/q_0}$ is a principal $m=N/\alpha$ root.
Indeed for $s=\alpha/q_0$ and $0<j<m$, the old $K=sm$ character sum
at exponent $sj$ is $s$ copies of the $m$-sum. Odd $s$ is one in
characteristic two, so that sum vanishes.

Append unused prime factors until $\alpha'=\alpha J$ satisfies
\[
 D'=\varphi(\alpha')\ge(R_0+L_0+5)^{64}.
\]
Only $O(\ell)$ factors are needed because each multiplies $D$ by more
than $L_0^2$. There are $\Omega(L_0/\ell^2)$ available factors.
The first-crossing choice implies
\[
 \log\alpha'=O(\ell^4),\qquad
 \varphi(\alpha')/\alpha'>1-1/L_0,
 \quad m'D'=\Theta(N),\quad m'=N/\alpha'.
\]
We do not claim the enlarged divisor meets the original definition's
literal numerical upper cutoff; these displayed estimates are sufficient.
We retain the stronger density estimate as well. The original divisor
has $O(\ell^2)$ prime factors, and enlargement adds only $O(\ell)$,
all greater than $L_0^3$. Thus
\begin{equation}
 \label{eq:strong-density}
 \theta:=1-\frac{D'}{\alpha'}
   =1-\prod_{p\mid\alpha'}(1-p^{-1})
   \le\sum_{p\mid\alpha'}p^{-1}
   =O(\ell^2/L_0^3).
\end{equation}
This follows from the actual divisor construction in
\cite[Proposition 3.8, proof]{cyclomult}, rather than just its
weaker admissible-divisor definition.

The substitution $Y\mapsto U^J$ defines a unital homomorphism
\[
 \F_2[Y]/\Phi_\alpha(Y)\longrightarrow
 S=\F_2[U]/\Phi_{\alpha'}(U),
\]
because $\Phi_{\alpha'}(U)$ divides $\Phi_\alpha(U^J)$.
Write $\iota$ for this homomorphism. It preserves principal-root
character sums. The required root is
$\omega'=\iota(\omega)^J$: it satisfies $(\omega')^{m'}=1$.
For $0<j<m'$, split the old $m=Jm'$ terms into $J$ equal periods
to obtain
\[
 0=\sum_{a=0}^{m-1}\iota(\omega)^{Jja}
   =J\sum_{b=0}^{m'-1}(\omega')^{jb}.
\]
Odd $J$ is one in characteristic two, proving principality.
Thus no new root oracle is needed.

Pad the remaining prime axes by $n_i=p_i+1$. Then
$M/m'=\prod_i(1+1/p_i)=1+O(L_0^{-2})$.
Since $\alpha'>L_0^{64}$, for sufficiently large $N$,
$L=\lceil\log_2(2M)\rceil\le L_0$ and $L=\Theta(L_0)$.
The native radix hypotheses and $D'\ge P_1(R_0+1,L)$ hold together.
Setup polynomial in $R_0,D',L_0$ is $N^{o(1)}$ and is charged once.

\section{The coupled CRT index permutation}

Write $m'=\prod_{i=1}^t p_i$, with least-significant-first digits
\[
 x=a_1+p_1a_2+\cdots+(p_1\cdots p_{t-1})a_t.
\]
Its physical lexicographic axis order is $[p_t]\cdots[p_1]$, already
ordinary integer index order. In one lexicographic scan embed these
records in the $[p_i+1]$ grid, preserving their order, and zero-fill
every address containing at least one dummy coordinate.

\begin{lemma}[Parallel CRT tree]
\label{lem:crt}
The permutation $a\mapsto(x\bmod p_i)_i$ and its inverse have cost
$O(m'D'L_0^\gamma\polylog L_0)$.
\end{lemma}
\begin{proof}
For a contiguous block $I=A\cup B$ with $A$ lower significance,
write $x_I=x_A+P_Ax_B$, $P_A=\prod_{i\in A}p_i$.
After converting both children to their local residues, the left
residues are already correct for $x_I$. Each right target becomes
\[
 z_i\longmapsto P_Az_i+x_A\pmod{p_i}.
\]
This is a permutation because $P_A$ is a unit modulo $p_i$.
The left residue address determines $x_A\in[0,P_A)$ by ordinary
integer CRT arithmetic. Fix the target dummy, and use the identity
when a relevant left control is dummy. This is an explicit bijection
on every fiber and preserves every dummy-pattern class at completion.

At each level of a balanced tree take the union of all right axes as
$T$, and all left/passive axes as $C$. The disjoint node merges are
one controlled tensor from Lemma~\ref{lem:controlled}. There are
$O(\log(t+1))$ levels. This schedule avoids two separate full-volume
child executions per node. Reverse the stages and affine maps for
the inverse.
\end{proof}

\section{Characteristic-two local Fourier compilation}

For an odd prime $p$, let $u\in S$ be a principal $p$-root. Every field
component has exact order $p$, so $1-u^j$ is a unit for $0<j<p$.
Set $H_j=\prod_{a=1}^j(1-u^a)$ and $H_0=1$. The Newton matrix
\[
 N_{j,k}=\prod_{a=0}^{k-1}(u^j-u^a),\quad j\ge k,
\]
is diagonal--lower-Toeplitz--diagonal, with
$N_{j,k}=H_jH_{j-k}^{-1}(-1)^k u^{k(k-1)/2}$.
The symmetric Newton factorization is
$F_p(u)=N\operatorname{diag}(N_{kk})^{-1}N^T$.
All pivots remain units over the finite etale algebra.

The explicit formal triadic circuit in
Lemma~\ref{lem:formal-convolution}, with one operand fixed, supplies a
linear convolution DAG of size $O(b\log^8b)$ and depth $O(\log^2b)$
over every commutative $\F_2$-algebra. Bounded-degree copy trees are
included. The local appendix also gives explicit constant/linear
gate rules for this formal specialization, with no Boolean
idempotence identity substituted in the coefficient ring.

The borrowed-coordinate mechanism of \cite{local} restores arbitrary
dirty workspace by replaying once with the true sources and once with
the sources suppressed. Its algebra works in characteristic two.
Inside a width-$v$ Toeplitz recursion, choose chunk size
$b=c v/\log^8v$, with $c$ a sufficiently small fixed constant. The
two chunks and all gate coordinates fit within the same $v$ positions.
There are $O(\log^{16}v)$ source/target chunk pairs, each using
$O(\log^2v)$ replay rounds. A balanced recursion gives
$O(\log^{19}p)$ rounds on exactly $p$ coordinates, as proved in
Lemma~\ref{lem:exact-width-toeplitz}. Its source, target and borrowed
ports are disjoint; output-copy wires are counted inside the DAG.
All borrowed values
return before the next chunk pair. The finite small range is direct.

Work in $B=S\otimes_{\F_2}\F_4$ and take central
$\eta^2+\eta+1=0$. Every coefficient $\mu\in B$ is the sum of two
units: let $e$ be one on the components where $\mu=1$ and zero
elsewhere, and put $d=1+e(\eta+1)$. Then $d$ and $\mu+d$ are units.
All component degrees equal $r_B=\operatorname{lcm}(\operatorname{ord}_{\alpha'}2,2)$;
thus $e=1+(\mu+1)^{2^{r_B}-1}$ can be prepared by polynomially many
small-algebra operations. For unit $t$,
\[
 U(t)=\begin{pmatrix}1&t\\0&1\end{pmatrix}
 =\operatorname{diag}(t,1)HJ\operatorname{diag}(t^{-1},1),
 \qquad U(\mu)=U(d)U(\mu+d).
\]
Here is one uniform slot word for the whole round. Pad each $p$ axis
by its one dummy and complete a partial shear matching by pairing
its unused coordinates with coefficient $\mu=0$. Order every used
pair as (target,source), and send these pairs to canonical adjacent
pairs by a local permutation. Set $u_0=d$, $u_1=\mu+d$ and
$D(t)=\operatorname{diag}(t,1)$. In rightmost-first matrix order,
\[
 U(\mu)=D(u_0)HJ D(u_0^{-1}u_1)HJ D(u_1^{-1}).
\]
All diagonal entries are units. An unused pair has $u_0=u_1=1$,
so its completed word is $(HJ)^2=I$; it need not be fixed at an
intermediate letter. Tensoring uses exactly two full canonical $H$
slots, two full $J$ slots and three axis-factorized diagonal slots.
The pair relabeling and its inverse each cost one generic tensor
permutation from \eqref{eq:generic}; charge the two $J$ slots there
as well. Thus a round has two native-$H$ calls and at most four
generic permutation calls. No separate selection of $H$ versus $I$
is needed on this local path. Each completed local axis map is
$F_{p_i}\oplus[1]$ on its native-valid range $n_i=p_i+1$.
The tensor preserves each per-axis dummy-pattern subspace; it need
not fix every full address containing a dummy, since other axes still
transform. Consequently an initially zero-filled dummy complement
stays zero at the completed tensor boundary. Keep the full padded
grid through these operations, including every intermediate letter.

\begin{lemma}[Independent runtime word arithmetic]
\label{lem:old-word}
There is a uniform fixed-tape algorithm for ordinary binary polynomial
products at length $v$ in $O(M_{\rm old}(v))$ time, with bound function
\[
 M_{\rm old}(v):=v\log^2(2v)\log\log(4v).
\]
It uses classical integer multiplication, independently of the proposed
main construction and of the formal convolution circuit supplier.
\end{lemma}
\begin{proof}
For $v\ge2$ set $k=\lceil\log_2(v+1)\rceil$ and lift
$\sum_{i<v}a_iX^i$ to $\sum_{i<v}a_i2^{ki}$, with $a_i\in\{0,1\}$.
Every integer convolution coefficient is at most $v<2^k$.
Consequently the product's $k$-bit chunks are the exact integer
coefficients, without carries between chunks. Their low bits give the
ordinary $\F_2$ product. The encoding, zero padding, chunk extraction
and optional bit-order reversals are sequential scans costing $O(vk)$.
Counter resets are charged to their ensuing chunks. The classical
integer algorithm of \cite{schonhage-strassen} takes
$O(B\log(2B)\log\log(4B))$ steps on a fixed multitape machine for
$B=vk$ bits, proving the bound. Marking visited work-tape intervals
allows head returns and cleanup within the same order of time.
Length one is direct. The bound function is stable under fixed-factor
length changes.
\end{proof}

Tensor scalar coefficients are prepared by prefix expansion in the
current native order, as detailed in the local appendix. The total
number of products is below $2M$, not one product per axis at every
record. A $B$
symbol has $2D'$ bits. Each $B$ multiplication or fixed-scalar
application costs $O(M_{\rm old}(2D'))=O(M_{\rm old}(D'))
=D'\polylog L_0$ bit steps using Lemma~\ref{lem:old-word} and the
explicit fixed-modulus reduction in the local appendix.
In contrast, the routing compiler's $\F_4$
scalars are fixed two-bit maps costing $O(D')$. These two costs must
not be conflated. Project back to $S$ only after the complete Fourier
tensor has returned to that subspace.
Each synchronized replay round has a fixed number of calls at the
generic/native $\gamma$ tariff and a fixed number of these scalar
stages. There are $O(\log^{19} R)=\operatorname{polylog}L_0$ rounds,
so the local Fourier tensor retains the same $\gamma$ exponent.

\section{Ordinary multiplication and the time ledger}

Use Lemma~\ref{lem:host-initialization} to construct
$N=\Theta(n)>2n-2$ and its certified factor list, then enlarge the
divisor as above. Initial input scans cost $O(n)$ and all host
initialization is included in the $N^{o(1)}$ setup term below.
The top Agarwal--Cooley isomorphism from \cite[Lemma 2.2]{cyclomult} is
\[
 X\longmapsto Y^cZ,\quad c=(m')^{-1}\bmod\alpha',
\]
from $\F_2[X]/(X^N-1)$ to
$\F_2[Y,Z]/(Y^{\alpha'}-1,Z^{m'}-1)$.
Its coprime factors $\alpha',m'$ both exceed two for sufficiently large
$n$: $\alpha'$ tends to infinity and $\log\alpha'=o(\log N)$.
It uses one classical rectangle transpose and cyclic row shifts in
$O(N\log\alpha')=N\polylog L_0$ bit time, in either direction.
Explicitly, an exponent $k=r+m'j$ goes to
$(Z\text{-index},Y\text{-index})=(r,cr+j\bmod\alpha')$;
this fixes the row-shift convention and its inverse.
Since $\alpha'$ is odd, $\Phi_{\alpha'}$ is coprime to
$\psi=(Y^{\alpha'}-1)/\Phi_{\alpha'}$.
Put $f=\Phi_{\alpha'}$ and $d=\deg\psi=\alpha'-D'$.
If $d=0$ the coefficient split is the identity. Otherwise use
Lemma~\ref{lem:fixed-monic-remainder} on each $\alpha'$-bit row
for both projections, including the possibly much shorter modulus
$\psi$. The latter's prepared reciprocal has length $\alpha'-d=D'$.
Choose fixed Bezout polynomials $uf+v\psi=1$, with
$\deg u<d$ and $\deg v<D'$. The idempotents
$e_f=v\psi$ and $e_\psi=uf$ have degree $<\alpha'$, and reconstruct
\[
 A=(A_f e_f+A_\psi e_\psi)\bmod(Y^{\alpha'}-1).
\]
Two padded length-$\alpha'$ word products and a linear cyclic fold
suffice for this reconstruction. Projections and reconstruction thus
  cost $O(M_{\rm old}(\alpha'))$ per row, and
  \[
    O(m'M_{\rm old}(\alpha'))=N\polylog L_0
  \]
  overall. All reciprocals,
Bezout coefficients and idempotents are public polynomial-$\alpha'$
setup, charged once. The two projected streams have total length
$m'(D'+d)=N$ bits.

For the large component let $\omega$ be the principal $m'$ root and
\[
 E_i=(m'/p_i)((m'/p_i)^{-1}\bmod p_i),\quad \omega_i=\omega^{E_i}.
\]
The CRT idempotents give
$\omega^{xy}=\prod_i\omega_i^{(x\bmod p_i)(y\bmod p_i)}$.
Lemma~\ref{lem:crt} and the local Fourier tensor therefore compute
the true ordinary cyclic DFT. Two forward transforms, point products
and one inverse compute the product in $S[Z]/(Z^{m'}-1)$.
Retain the full native grid through these completed operations;
its dummy complement remains zero at each such boundary and under
the point product, so it can be cropped after inverse routing.
The inverse uses inverse roots and reversed index routing; its odd
normalization is one in characteristic two.

Here the large coefficient ring obtained by projection is
$\F_2[Y]/\Phi_{\alpha'}(Y)$, identified with $S$ by renaming $Y$ to $U$.
This identification is distinct from the earlier setup homomorphism
$Y_{\rm old}\mapsto U^J$ from the smaller ring
$\F_2[Y_{\rm old}]/\Phi_\alpha(Y_{\rm old})$. That homomorphism constructs
the root; it is not applied a second time to the new input words.
In this section $\omega$ denotes the already constructed $\omega'$.
Indeed if $J>1$, $U^J$ has order $\alpha$, rather than $\alpha'$,
in every field component of $S$. Hence $\Phi_{\alpha'}(U^J)$ is a
unit, so that repeated substitution would not even descend from the
new quotient.

Frequency order needs its own verification. If local roots instead have
exponents $b_i h_i$, with $b_i=m'/p_i$ and $h_i$ a unit modulo $p_i$,
an output digit $s_i$ describes frequency residue $b_i h_i s_i$.
The idempotent choice $h_i=b_i^{-1}$ makes that residue $s_i$.
For a coordinate permutation $Q$, however, $A=QF$ still satisfies
$A^{-1}((Aa)\odot(Ab))=F^{-1}((Fa)\odot(Fb))$.
Thus matching inverse and product checks alone do not certify a
literal Fourier index convention.

The complement has $d=\deg\psi=\alpha'-D'$ and
$v=dm'=N\theta=O(N\ell^2/L_0^3)$ by \eqref{eq:strong-density}.
If $d=0$, omit this component.
Pack $Y^jZ^i\mapsto T^{2di+j}$ in $\F_2[T]$. There is no overlap
between coefficient-product slots because $j+j'<2d$.
Each input has degree at most $2v-d-1$ and the product degree at most
$4v-2d-2$. Use Lemma~\ref{lem:old-word} at input length $2v$.
Crop its known zero tail to $2d(2m'-1)$ coefficient slots and unpack
$2m'-1$ rows of $2d$ bits, each with final bit zero.
Reduce their degree-$\le2d-2$ polynomials modulo the fixed monic
$\psi$ using Lemma~\ref{lem:fixed-monic-remainder}, then fold
$Z^{m'}=1$. All packing, unpacking and folding scans cost $O(v)$;
when $d=1$ the row remainders are direct.
The total cost is
$O(M_{\rm old}(2v)+m'M_{\rm old}(d))
 =O(M_{\rm old}(v)+m'M_{\rm old}(d))$. Since
$d\le v\le N$, each logarithmic factor is bounded using $L_0$,
so the total is
\[
 O(vL_0^2\ell)=O(N\ell^3/L_0)=o(N).
\]
The scans and all reductions are included. The weaker estimate
$v<N/L_0$ alone would not pay for this integer-packing fallback;
the stronger density is essential here. No integer packing of the
full ordinary input is used.

\begin{lemma}[Typed top-level execution]
\label{lem:typed-execution}
The preceding component algorithms compose on a fixed number of streams,
with ordinary coefficient input and output. Their complete calls use
the stated tariffs; no coefficient projection or outer native-dummy crop
occurs inside a Fourier word.
\end{lemma}
\begin{proof}
Use the following order, with $a,b$ denoting the two inputs.
\begin{enumerate}
\item Count and copy the coefficient lists, construct the host and all
length-dependent tables, and zero-fill both lists to $N$ coefficients.
The live top-level input streams contain $2N$ bits. Their cost is
$O(n)+N^{o(1)}$ before the linear zero-fill.
\item Apply the top isomorphism to each list, producing $m'$ consecutive
$\alpha'$-bit coefficient rows. Project each row to the $f$ and $\psi$
components. The two large streams contain $m'$ words of $D'$ bits;
the two small streams contain $m'$ words of $d$ bits. Park the small
streams on a fixed separate pair of tapes. The maps and projections cost
$N\polylog L_0$. The large words are identified with $S$ by the variable
rename just specified.
\item Embed each large word in $B=S[\eta]/(\eta^2+\eta+1)$ as $a+0\eta$,
and zero-fill its native dummy grid. Keep these two $M$-word $B$ arrays
through both forward transforms, their point product and the inverse.
They occupy $4MD'=O(N)$ bits at these top-level boundaries.
Transform one while the other waits on its own tape; constant whole-array
passing and parking scans cost $O(N)$. No callee traverses the parked
array.
\item To specify every completed transform, let $E:B^{m'}\to B^M$
be the zero-fill embedding from integer mixed-radix order. Let $Q$ be
the completed extended CRT permutation, and put
\[
 \mathcal T_{\omega}=\bigotimes_i(F_{p_i}(\omega^{E_i})\oplus[1]),
 \qquad \mathcal D_{\omega}=Q^{-1}\mathcal T_{\omega}Q.
\]
The CRT exponent identity gives
\begin{equation}
 \label{eq:typed-dft}
 \mathcal D_{\omega}E=E F_{m'}(\omega).
\end{equation}
Thus the output is again in ordinary integer frequency order, in the
same native grid. The whole operators $Q$ and $\mathcal T_\omega$
preserve every dummy-pattern subspace. Their intermediate words may
mix the auxiliary coefficient lanes or dummy values; only completed
boundaries are used in this identity. Their valid block has entries
in $S$, so a complete transform returns $S$-valued words when its
input was $S$-valued. Keep the unused $\eta$ lane allocated nonetheless.
\item Replace the first completed transform array by the coordinatewise
$B$ product of the two arrays, scanning them together, and erase the
second array. On valid coordinates this is exactly the $S$ product;
on every dummy coordinate it is zero. Performing all $M$ products,
without an exceptional-address lookup, costs
$O(MM_{\rm old}(D'))=N\polylog L_0$. Apply
$\mathcal D_{\omega^{-1}}$ to the first array. The principal-root
character sums imply
$F_{m'}(\omega^{-1})F_{m'}(\omega)=m'I=I$.
Equation~\eqref{eq:typed-dft} and the convolution identity therefore
give the cyclic product in $S[Z]/(Z^{m'}-1)$.
Explicitly, for the projected coefficient vectors $a_f,b_f$ the
$j$th inverse coefficient is
\[
 \sum_{r,s=0}^{m'-1}(a_f)_r(b_f)_s
       \sum_{k=0}^{m'-1}\omega^{k(r+s-j)}
   =\sum_{r+s\equiv j\pmod {m'}}(a_f)_r(b_f)_s.
\]
The inner character sum is zero off the stated congruence and is
$m'=1$ on it. Thus the complete retained-grid computation is
$\mathcal D_{\omega^{-1}}
 ((\mathcal D_\omega E a_f)\odot(\mathcal D_\omega E b_f))
 =E(a_f*_{\rm cyc} b_f)$.
Only now remove the zero $\eta$ lane and crop the zero dummy complement,
using monotone native-address scans costing $O(N)$. The result is
$m'$ consecutive $D'$-bit words in ordinary $Z$ order.
Generate the dummy test once per record using the fixed degree-twelve
address procedure; its cost is paid by $D'\ge P_1$.
\item Compute the parked small-component product by the independent
packing and fixed-monic reductions above, at cost $o(N)$, obtaining
$m'$ consecutive $d$-bit words, omitting this branch if $d=0$.
Reconstruct the $\alpha'$-bit rows from
the two components and apply the inverse top isomorphism at cost
$N\polylog L_0$. Extract the first $2n-1$ coefficients in a final
$O(N)$ scan. All other top-level streams are erased and their heads
returned by scans charged to their lengths.
\end{enumerate}
There are three complete transforms, each with a constant number of CRT
maps and the stated synchronized local rounds. Their total is
$O(NL_0^\gamma\polylog L_0)$, since their $B$ payload width is $2D'$
and hence satisfies the same metadata lower bound. The number of
top-level streams is fixed; setup archives and recursive scratch use
the separately charged tape families, rather than one fresh tape per
axis, table or recursion level.
\end{proof}

The full-volume routing and Fourier arrays have $V=\Theta(N)$ and
address width $L=\Theta(L_0)$. The smaller complement arrays have
the separate budget above. The stage counts are fixed powers of $\log L_0$.
The complete bound furnished by the construction is
\[
 O\bigl(NL_0^\gamma(\log L_0)^C\bigr)
   +N\polylog L_0+N^{o(1)}.
\]
With $\epsilon=2^{-36}$, the reconstructed interchange and phase
certificates give $\tau=1-\epsilon$ and $\sigma<\tau$. Hence
\[
 \gamma=\frac1{1+\epsilon},\qquad
 1-\gamma=\frac{\epsilon}{1+\epsilon}>2^{-37}.
\]
For fixed $C$, all secondary logarithms are therefore absorbed in
$O(n(\log n)^{1-2^{-38}})$. The same exponent is retained through
the controlled compiler and CRT tree; their extra work is linear
metadata and fixed powers of $\log L_0$, not a larger power of $L_0$.
The effective finite initial range uses a fixed direct branch.
Reconstruction and inverse top routing return the ordinary cyclic
$N$ coefficients; degree $2n-2<N$ means the first $2n-1$ coefficients
are exactly the ordinary dense polynomial product.

\paragraph{Order of choices and one cutoff.}
The finite motifs, avoiding prime, generator degrees and exponents are
fixed first. Host construction then fixes $N,L_0,\ell,R_0$ before
enlargement. The threshold for $D'$ uses $L_0$, and the resulting
native width satisfies $L\le L_0$; it is not fed back into the host.
Only $O(\ell)$ extra factors are consumed from
$\Omega(L_0/\ell^2)$ available factors. All setup degrees and secondary
logarithmic powers are fixed constants of these algorithms.
The imported host constants are effective \cite[Section 1]{cyclomult};
the other constants come from the printed finite generators and
constructive bounds. One maximum of their finitely many eventual
thresholds therefore supplies the direct-branch cutoff.
For explicit final logarithm absorption, write $\delta=2^{-38}$ and
bound the main term by $ANL_0^\gamma(\ln L_0)^C$ with fixed $A\ge1$.
For
\[
 t=\ln L_0\ge\max\{1,(2C/\delta)^2,2\ln(A)/\delta\},
\]
$\ln t\le\sqrt t$ gives $A(\ln L_0)^C\le L_0^\delta$.
Since $1-\gamma>2\delta$, this fits the stated $1-\delta$ exponent.
No input-dependent exponent or length advice is selected.

\begin{theorem}[Main theorem]
\label{thm:main}
There is one uniform deterministic machine with a fixed finite alphabet
and a fixed finite number of one-dimensional tapes which multiplies two
ordinary dense degree-$<n$ polynomials over $\F_2$, for every $n\ge1$,
and returns their $2n-1$ ordinary coefficients in time
\[
 O\!\left(n(\max\{1,\lceil\log_2 n\rceil\})^{1-2^{-38}}\right).
\]
Its construction uses no unproved prime-distribution hypothesis.
\end{theorem}
\begin{proof}
Lemma~\ref{lem:host-initialization} supplies the actual admissible length,
factor list, divisor and root from $n$ in charged $N^{o(1)}$ setup,
with $N=\Theta(n)>2n-2$. Word enlargement keeps $m'D'=\Theta(N)$ and makes
all controlled-routing coefficients computable within their word cost.
The reconstructed interchange, finite phase and native endpoint
appendices prove Proposition~\ref{prop:suppliers}; the ragged and
controlled compilers therefore price the actual index permutations.
Lemma~\ref{lem:crt} supplies the ordinary-to-prime-residue index map.
The local appendix gives the exact prime Fourier words over the
specified finite etale algebra, with separately charged classical scalar
maps. The preceding component calculation computes the large and
complementary coefficient-ring products, reconstructs the cyclic
$N$-product, and extracts the ordinary product. The displayed time
ledger absorbs every fixed secondary logarithm in the stated exponent.
Lemma~\ref{lem:typed-execution} specifies the actual stream composition,
root/coefficient types and the single final auxiliary-lane/dummy crop.
The finite tape and pad composition below applies to these actual
subroutine calls. A direct finite branch handles the bounded initial
range, with one enlarged uniform constant.
\end{proof}

\begin{corollary}[Each fixed finite extension of $\F_2$]
For every fixed integer $s\ge 1$, multiplication of ordinary dense
degree-$<n$ polynomials over $\F_{2^s}$ can be performed by one
deterministic machine with a fixed finite alphabet and fixed finite
number of tapes, in running time
\[
 O_s\!\left(n(\max\{1,\lceil\log_2 n\rceil\})^{1-2^{-38}}\right).
\]
This is a derived corollary of Theorem~\ref{thm:main}; its constants and
machine may depend on $s$.
\end{corollary}
\begin{proof}
Fix a basis of $\F_{2^s}$ over $\F_2$ and write each input coefficient
in that basis. By bilinearity, each output coordinate of the product is
a fixed $\F_2$-linear combination of the $s^2$ ordinary binary
coefficient products of the input coordinate polynomials. Run
Theorem~\ref{thm:main} on those $s^2$ pairs in sequence, reusing the
same finite tape family, and combine the resulting coefficient lists
using the fixed multiplication table of the chosen basis. The $s^2$
calls cost $O_s(n(\log(n+2))^{1-2^{-38}})$; the coefficientwise
recombination costs $O_s(n)$. For each fixed $s$ this gives the stated
bound. No uniformity as $s$ grows with $n$ is asserted.
\end{proof}

\paragraph{Tape and padding composition.}
Host I/O, classical integer SS, ring arithmetic, metadata generation, generic routing,
H2 roles, weighted tiles and SWAP have disjoint fixed tape families.
Each recursive family reuses its tapes with current-parent stack
parking. No child traverses parked ancestors. The caller waits while
a subroutine operates on its actual stream; every passing/return scan
is included in that subroutine's cost. The native dummy grid is kept
through complete routing/Fourier operations. Private rectangle, tile,
role and suffix pads have the precise lifetimes in the machine ledger;
none is repeatedly multiplied into recursive child mass.
The complete family call graph, head-return rules and crop table are
specified in the included machine-composition record
\path{supporting/STRICT_BINARY_MACHINE_COMPOSITION_ROOT_20261030.txt},
Sections~III--V.

\paragraph{Verification evidence and remaining review.}
A previously recorded ragged-network witness checked 42,090 permutations at even
sizes at most 12. A previously recorded guarded C++ witness checked 768 cases on
native axes $(4,6,8)$, including all 105 valid CRT addresses for primes
$(3,5,7)$, inverse maps and exact dummy patterns. Omitting the control
gauge fails. Both install 64 MiB/10 CPU-second limits before algebra.
They validate finite identities and conventions only. The bundled
machine-composition derivation and proof map document the integration
dependencies; internal audits do not replace external mathematical review
or establish novelty. These existing witness reports were not rerun for
this release copy.
A previously recorded guarded fixed-array C++ witness checked all $65\,536$
coefficients of the small product algebra
$(\F_2[Y]/\Phi_{15})\otimes\F_4$ for the selector/unit split and
the local two-$H$ word, with independent bilinear multiplication checks.
Previously recorded mixed-radix tensor fixtures checked scalar-prefix expansion,
arbitrary incoming values and dummy restoration. Wrong middle factors
and stale scalar alignment are rejected. This remains bounded algebraic
evidence, not a test of the asymptotic tape bound.
A previously recorded bounded whole-composition witness used $N=255=17\cdot15$,
both coefficient CRT components, the affine residue-index tree,
compiled $3$- and $5$-point Newton words and ordinary reconstruction.
It checks the literal forward/inverse DFT on all $15$ basis vectors
and $96$ ordinary binary products through degree bound $n=128$.
Omitting row shifts, the complement or the required extension lifetime
is detected. Consistently changing primitive roots preserves its
roundtrips/products while failing the literal Fourier comparison.
The witness uses direct finite primitive branches and does not test
the growing-size running-time theorem.
A previously recorded guarded C++ witness checked the elementary triadic product
over $\F_2[\epsilon]/(\epsilon^2)$: $130$ ordinary products through
length $512$, a detected wrong-overlap control and a nilpotent-law
check. It uses direct finite DFTs and base cutoff $32$ to exercise
the quotient-product recursion. The proof's radix-three depth and
analytic cutoff $6^6$ are separate; this witness establishes neither
the asymptotic circuit bounds nor the main theorem.
An expanded source-hashed finite check also verified the 776-letter
$E$ word, the explicit (A.17) conjugation and monomial scalar ordering.
Ordinary-squaring and stale-scalar negative controls are detected.
Its fixed four-by-four field and 16-bit binary maps do not test the
growing-size native routing or time bound.

\appendix
\section{Finite ternary phase router and a positive binary $H_2$ word}
\label{app:finite-phase}

The scalar bank-exchange and frame-label architecture follows the
finite-network method of~\cite{swap}; its characteristic-two and
ternary-phase adaptation is made explicit below. These finite-network
antecedents are credited independently of the present tape-time claim.

This appendix records the finite algebraic and rank-saving part of the
binary $H_2$ supplier. It distinguishes three coefficient/address domains:
the address-frame geometry is over $\mathbb F_3$, the phase kernel and
dirty-selector payload arithmetic are over $\mathbb F_4$, and the fixed
positive word is over its codec field $\mathbb F_{16}$. The tape costs for
the equal-chunk interchange and weighted native mask are separate machine
suppliers; no field element in these algebraic arguments is treated as a
unit-cost operation on a growing coefficient word.

\subsection{The ternary phase kernel}

Let $\eta\in\mathbb F_4$ satisfy $\eta^2+\eta+1=0$. On $\mathbb F_4^3$ put
\[
 C=\eta I_3+\eta^2J_3
 =\begin{pmatrix}1&\eta^2&\eta^2\\
 \eta^2&1&\eta^2\\\eta^2&\eta^2&1\end{pmatrix}.
\]
Then $C^3=I$. Let $\chi(a)=\eta^a$ for $a\in\mathbb F_3$, and let
\[
 (\mathcal F_3)_{x,z}=\chi(xz),\qquad x,z\in\mathbb F_3.
\]
Since $1+\eta^a+\eta^{2a}=0$ for nonzero $a\in\mathbb F_3$,
$\mathcal F_3^2$ is the coordinate-negation permutation. Thus
$\mathcal F_3^4=I$ and $\mathcal F_3^{-1}=\mathcal F_3^3$ (equivalently,
$(\mathcal F_3^{-1})_{x,z}=\eta^{-xz}$).
Direct multiplication gives
\[
 C=\mathcal F_3\,\operatorname{diag}(1,\eta,\eta)\,\mathcal F_3^{-1}.
\]
Thus the eigenphase of $C$ at frequency $z$ is $\eta^{z^2}$. The inverse
is $C^2$ and has eigenphase $\eta^{-z^2}$.

For an address space with a nondegenerate symmetric form over $\mathbb F_3$,
write $T_v f(x)=f(x+v)$. If $v$ has norm $\nu=\langle v,v\rangle\in
\{1,2\}$, define
\[
 K_v=I+\eta^{-\nu}(T_v+T_{-v}).                         \tag{A.1}
\]
On a Fourier character of frequency $z$, the two translations have
eigenvalues $\eta^{\langle v,z\rangle}$ and
$\eta^{-\langle v,z\rangle}$. For $a=\langle v,z\rangle\in\mathbb F_3$,
\[
 1+\eta^{-\nu}(\eta^a+\eta^{-a})=\eta^{a^2/\nu}.     \tag{A.2}
\]
Indeed both sides are $1$ for $a=0$; for $a=\pm1$, checking
$(\nu,a^2)=(1,1)$ and $(2,1)$ gives $\eta$ and $\eta^2$ respectively.
Consequently $K_v$ is precisely the one-direction quadratic phase for
the line $\langle v\rangle$. Orthogonal directions commute, and an
orthogonal direct sum of directions has the product of their phases.
Replacing $v$ by $-v$ does not change the line; norm $2$ replaces $C$ by
$C^{-1}$.

For a nondegenerate subspace $U$ with an orthogonal basis
$v_1,\ldots,v_r$, define $C_U=\prod_jK_{v_j}$. In Fourier coordinates
this is the multiplier
\[
 (\mathcal F C_U\mathcal F^{-1})_{z,z}
   =\eta^{\sum_j\langle v_j,z\rangle^2/\langle v_j,v_j\rangle}.
                                                               \tag{A.3a}
\]
Thus orthogonal direct sums multiply their $C_U$ operators. If
$U\subseteq V$ and $V=U\perp W$, then $C_V=C_UC_W$; a rank-$r$
edge between such frames needs exactly $r$ one-direction children.
The phase frame on a role is a proof invariant: the same $\mathbb F_4$-linear
phase operator is applied to aligned wire arrays, so it commutes with
pointwise XOR/addition. It is a linear operator on payload arrays, not an
address relabeling; this is distinct from the address permutations used
in the SWAP supplier.
At completion the role permutation pairs each input frame with its
specified output frame. Choosing the X-source inverse line frame makes
the endpoint product the full phase, while its first edge has difference
$P_U-(-P_U)=2P_U$ and therefore one rank-one child. This is the source
correction counted below.

\subsection{The $48$-triple frame over $\mathbb F_3$}

Fix $H=48$, let $\mathcal T=\binom{[H]}3$, and join distinct $S,T$ when
$|S\cap T|=1$. Put
\[
 v=\binom{48}{3}=17296,\quad
 z=3\binom{45}{2}=2970,\quad N=v^3,\quad I=3v^2.
\]
The factor space is $F=\mathbb F_3^{49}$ with form
\[
 B=\operatorname{diag}(-I_{48},1).
\]
For $S\in\mathcal T$ use $t_S=(1_S,1)$. Then
\[
 \langle t_S,t_S\rangle=-3+1=-2=1\pmod 3,
 \qquad \langle t_S,t_T\rangle=1-|S\cap T|.
\]
Thus the line $L_S=\langle t_S\rangle$ is unit norm and adjacent triples
give orthogonal lines. The form is nondegenerate. To identify it with
standard dot-product coordinates, split the first $48$ coordinates into
$24$ disjoint pairs, and on each pair use
\[
 (a,b)\longmapsto(2a+2b,\,2a+b).
\]
The two output vectors for the basis inputs are $(2,2)$ and $(2,1)$;
their norms are $2$ and $2$, and their dot product is $2\cdot2+2\cdot1
=0$ in $\mathbb F_3$. The map is invertible (determinant $1$ in $\mathbb F_3$), so this is
an isometry from $\operatorname{diag}(-1,-1)$ to the standard dot
product. The last coordinate already has norm one.

For a stage prefix tensor space $A=F^{\otimes(j-1)}$, let $P$ be its
prefix line, $B_0=P^\perp$, and let $L_t$ denote the current triple line.
Since $P$ and $L_t$ are nondegenerate,
\[
 A=B_0\perp P,\qquad F=L_t\perp L_t^\perp,
\]
and
\[
 (P\otimes L_t)^\perp
   =(B_0\otimes F)\perp(P\otimes L_t^\perp).        \tag{A.3}
\]
The neighbor condition supplies $L_s\perp L_t$ whenever a side role
contains the ordered neighbor $(s,t)$. These decompositions prove the
following eight gate labels are nested/comparable exactly where their
wire edges require:
\[
\begin{array}{c|l}
 t_0& B_0\otimes L_t\\
 t_1& B_0\otimes F\\
 t_2& (B_0\otimes F)\perp(P\otimes L_t)\\
 t_3& A\otimes F\\
 t_4& B_0\otimes F\\
 t_5& (B_0\otimes F)\perp(P\otimes L_t^\perp)\\
 t_6& A\otimes F\\
 t_7& A\otimes F.
\end{array}                                                   \tag{A.4}
\]
In stage two the physical wire order is the inverse of the forward row
schedule with logical source and target exchanged; this keeps the same
physical X/Y label convention in (A.4). In particular, its incoming
physical X label is $(A\otimes L_t)\otimes Q$ and its outgoing label
$(A\otimes F)\otimes Q$; its incoming physical Y label is
$(B_0\otimes L_t)\otimes Q$ and outgoing label
$(P\otimes L_t)^\perp\otimes Q$. The latter is (A.3) for the next
stage, with the suffix tensor line $Q$ restored. This checks the
interstage edges as well as the within-invocation inclusions.

Here is the wire-level interpretation of the table. A stage fixes the
other two triple indices of a data role in $\mathcal T^3$; all displayed
labels are tensored with its fixed suffix line $Q$. Name an ordered
neighbor pair by its physical Y and physical X triple indices in
that stage, renaming the ordered pairs in the inverse second stage.
The side wire
$a_{(S,T)}$ meets, in order, the Y-side gate for $S$, the X-side
gate for $T$, the Y-side gate for $S$, and the X-side gate for $T$.
Their labels are rows $t_0,t_2,t_5,t_7$. Its middle inclusion is
precisely $L_T\subset L_S^\perp$, supplied by the neighbor condition.
Each center wire meets $t_1,t_3,t_4,t_6$, whose labels alternate
$B_0\otimes F$, $A\otimes F$, $B_0\otimes F$, $A\otimes F$.
Data X wires meet $t_2,t_3,t_6,t_7$ and data Y wires meet
$t_0,t_1,t_4,t_5$. These descriptions check every gate incidence;
multi-input sums can be decomposed into the same common-label XOR
gates without adding frame edges of nonzero rank.
Side and center scratch wires have identity input phase and full-space
output phase, connecting their first/last displayed labels to zero/full
labels. Initial Y data have identity phase. Final X data have full-space
phase and final Y data have the complement of their triple tensor line.
The initial X inverse line is the sole exceptional source frame.
The physical scratch rows also require the full phase: they were
arbitrary rows of the original array. It is their logical scalar
scratch values that the routing motif restores before the frame is
interpreted at the output.

For a fixed eligible mask sector, let $P$ be this scalar role permutation
acting on aligned arrays, and let $D_{\rm in},D_{\rm out}$ be the
block-diagonal input/output phase frames. A frame change on an edge
from $F$ to $G$ applies $GF^{-1}$. Common-label XOR gates commute with
their shared phase, so the entire framed program telescopes to
$D_{\rm out}P D_{\rm in}^{-1}$. If input role $r$ is sent to $P(r)$,
the endpoint operator there is
$C_{{\rm out},P(r)}C_{{\rm in},r}^{-1}=C_{\rm full}$ for data
\emph{and} scratch roles. Consequently
\[
 D_{\rm out}P D_{\rm in}^{-1}
     =P(I_W\otimes C_{\rm full}).                       \tag{A.4a}
 \label{eq:full-role-phase}
\]
The inverse role merge removes $P$, leaving the full phase on every
arbitrary input role array. Eligibility is the same mask-only predicate
on every row role; in an ineligible sector every frame is identity,
so the scalar $P$ is likewise canceled by that merge. No scratch
payload is initialized to zero in this identity.

\subsection{Arbitrary-scratch scalar routing and the finite rank sum}

For one invocation over $\mathbb F_2$, let $x$ be its data source, $y$
its target, $a$ its side scratch vector and $c$ its center scratch vector.
Let $V,J,G,R$ be the incidence maps in the companion scalar-motif
certificate. The chronological eight rows are
\[
\begin{aligned}
 y&\gets y-Ja, & y&\gets y-Rc,\\
 a&\gets a+Vx, & c&\gets c+Gx,\\
 y&\gets y+Rc, & y&\gets y+Ja,\\
 c&\gets c-Gx, & a&\gets a-Vx.
\end{aligned}
\]
where in the first two rows $a,c$ are their entry values. In
characteristic two subtraction equals addition, but the signs make the
scratch cancellation visible. The final two rows restore the center and
side scratch. Equivalently,
for arbitrary initial scratch $(a_0,c_0)$ the total target increment is
\[
 -Ja_0-Rc_0+R(c_0+Gx)+J(a_0+Vx)=(RG+JV)x.             \tag{A.5}
\]
For triples, $G_i,T=1_{i\in T}$, $R_{S,i}=1_{i\in S}$,
$V_{(S,T),T'}=1_{T=T'}1_{|S\cap T|=1}$ and
$J_{S,(S,T)}=1_{|S\cap T|=1}$. If $S=T$, $RG=3=1$ in $\mathbb F_2$
and $JV=0$. If $S\ne T$ with intersection one, both terms are one
and cancel; for intersections zero or two both vanish. Hence
$RG+JV=I$. Formula (A.5) proves exact addition and restoration for every
scratch assignment.

Run the forward schedule X-to-Y, its inverse schedule Y-to-X, then the
forward schedule again. On arbitrary banks
\[
 (x,y)\mapsto(x,y+x)\mapsto(-y,y+x)\mapsto(-y,x),
                                                               \tag{A.6}
\]
so in characteristic two this exchanges the banks. Every invocation has
disjoint scratch roles; all are fixed by the completed role permutation.
There are $v$ roles per bank, $vz$ ordered side roles and $H$ center
roles per invocation, hence
\[
 W=2N+I(vz+48).                                           \tag{A.7}
\]

For a label of dimension $r$ over $\mathbb F_3$, a wire operation on
that label uses its quadratic phase and a source/target frame change of
rank equal to the relevant dimension difference. The only downward
label transitions in (A.4) are the $t_3\to t_4$ center transitions.
Each loses $F=49$ dimensions and there are $48$ center coordinates in
each of $I$ invocations. Their total downward loss is
\[
 L=I\cdot48\cdot49.                                    \tag{A.8}
\]
The signed dimension telescope over all routed wire endpoints is
$WM-2N$, where $M=49^3=117649$: the two banks contribute the two $N$
endpoint deficits. Turning signed variation into absolute variation
adds twice every downward loss, so before the source correction the sum
of edge ranks is $WM-2N+2L$.

This count can also be read from the complete physical wire paths.
Put $d=49$ and $a=d^{j-1}$ at stage $j$. With the suffix line
suppressed, the side and center dimensions are respectively
\[
 0,\ a-1,\ d(a-1)+1,\ da-1,\ da,\ M,
 \qquad
 0,\ d(a-1),\ da,\ d(a-1),\ da,\ M .
\]
The side inclusions follow from $L_T\subset L_S^\perp$, with
nondegenerate rank-$(d-2)$ residual. The center loses exactly
$d$ dimensions at $t_3\to t_4$ and otherwise increases.
Physical X enters the next stage as $A\otimes L_T\subset t_2(T)$;
physical Y enters as $B_0\otimes L_S=t_0(S)$.
Their complete ordinary paths have no decreases and have total
rank $M-1$ each, using a formal positive X input line for this count.
The inverse middle stage retains this physical order: its inverse
logical rows $7,\ldots,0$, after bank exchange, have the same physical
incidences as chronological rows $0,\ldots,7$. The label table itself
is not reversed. Consequently the total ordinary rank is
\[
 2N(M-1)+(W-2N)M+2I\cdot48d=WM-2N+2L.
\]

Choose the X-source phase frame to be the inverse line phase on $U_a$;
all other source frames are identity and gate/sink frames are the
positive phases of (A.4). The endpoint identity on each X-to-Y route is
$C_{U_a^\perp}C_{U_a}=C_{F^{\otimes3}}$. On each of the $N$ first X
edges, the source label is $-P_{U_a}$ while the first gate label is
$+P_{U_a}$, so their difference is $2P_{U_a}$. This has rank one and
its phase is $C_{U_a}^{-1}$, since scaling the quadratic form by
$2=-1$ in $\mathbb F_3$ inverts its phase. These are $N$ genuine extra
children, not a rank-zero equal-dimension edge. Therefore
\[
 s=WM-2N+2L+N=WM-N+2L.                                \tag{A.9}
\]

At $H=48$ direct cancellation gives
\[
 \frac WN=\frac{9633881}{1081},\qquad
 \frac LN=\frac{441}{1081},\qquad
 1-\frac{s}{WM}=\frac{199}{9633881\cdot117649}.
                                                               \tag{A.10}
\]
The integer comparison
\[
 199\,2^{36}>12\cdot9633881\cdot117649
\]
and $\log_2 M<17$, $\ln2<7/10$ imply $\ln M<12$ and
\[
 \frac{199}{9633881\cdot117649}>12\,2^{-36}
   >1-M^{-2^{-36}}.
\]
Thus $s/W<M^{1-2^{-36}}$. This is the explicit strict role saving used
below; no numerical matrix-rank test is involved.

\subsection{From the role word to all active sectors}

An $A$-digit has states $0,1,2$ active and state $3$ idle. Represent
$d_i=2X_i+Y_i$. For any active mask $S$, $A^{\otimes h}$ restricts to
$C^{\otimes |S|}$ on the active coordinates and identity on idle
coordinates. All $4^h$ addresses remain in the stream; no sector is
gathered out.

The mask-routing identities use dirty address fields. Let $T_Z$ apply
one fixed transposition of active symbols where dirty selector bit
$Z_i=1$. It preserves the active mask $S$, but arbitrary dirty $Z$
can make it change paired sources as well as targets. Source-value
predicates are recomputed only after the completed $T_{J\wedge Q}$
word below, which changes matched targets only. For a mask-only selector
$J(S)$ the chronological word
\[
 Z\mathrel{\oplus}=J(S),\quad T_Z,\quad
 Z\mathrel{\oplus}=J(S),\quad T_Z                         \tag{A.11}
\]
restores arbitrary initial $Z$ and equals $T_J$. With a second dirty
field $Q$, use $J(S)\wedge Q$ in (A.11). If $F_a(d)_i$ says the paired
source of target $i$ currently has trit $a$, then the completed
$T_{J\wedge Q}$ changes targets only and preserves all paired sources.
Consequently
\[
 Q\mathrel{\oplus}=F_a(d),\quad T_{J\wedge Q},\quad
 Q\mathrel{\oplus}=F_a(d),\quad T_{J\wedge Q}              \tag{A.12}
\]
restores arbitrary $Q,Z$ and implements the target transposition
conditioned on source value $a$. The source/target matching must be
disjoint; this is what makes the second $F_a$ equal to the first at the
completed inner-operation boundary. A third dirty field is reserved for
the weighted tile auxiliary. The supplied weighted-mask lemma implements
each frozen XOR in (A.11)--(A.12) with its separate auxiliary restored;
its exact sparse boundary repair and interchange cost are not reproved
here.

The two active transpositions have explicit bit circuits
\[
 (0\ 1):\ Y\mathrel{\oplus}=Z\wedge\neg X,
 \qquad
 (1\ 2):\ X\mathrel{\oplus}=Z\wedge Y;\quad
 Y\mathrel{\oplus}=Z\wedge X;\quad
 X\mathrel{\oplus}=Z\wedge Y.                           \tag{A.13}
\]
They are constant sequences of frozen-control XOR shears. Applying the
two transpositions in chronological order $(1\ 2)$ then $(0\ 1)$ adds
one to an active trit; reverse order subtracts one. Conditioning on
source trit and composing these shears realizes a matching SUM
\[
 \text{target}_i\gets\text{target}_i+\text{source}_i\pmod3.
                                                               \tag{A.14}
\]
Three SUMs plus a negation implement a coordinate swap. All dirty
selectors and the tile auxiliary return to their arbitrary incoming
values after each completed operation.

For a node acting on the first $e$ active ranks, restrict controls and
endpoints to the first
\[
 u=\min\{h,\lceil C_0(e+\log(L+2))\rceil\}
\]
physical positions, where $L\ge h$ is the ambient address-width bound
in the native implementation. If fewer than $e$ active positions occur there,
define the operation to be identity; otherwise all required ranks are
present. Restrict every raw $T_Z$ transposition and both selector-bank
toggles to this prefix, not just the completed selector $J$. Insert
inherited eligibility into $J(S)$; it is not a frozen control on an
inner bit shear whose target can change $S$ temporarily. Individual
inner shears still read only the other bit plane and the selector bank.
Their completed transposition preserves $S$. The comparison with the
full intended route is made only after a \emph{complete logical}
mask-selected transposition, source-conditioned operation (A.12), SUM
or recursive basis placement on the first $e$ active ranks, with all
its dirty selectors restored. Write this
prefix implementation as $P_u$ and its intended permutation as $T$.
On a good eligible prefix all intended endpoints are inside it, so
$P_u=T$. On a bad or ineligible prefix the selector $J$ is zero:
the two raw transpositions in (A.11) cancel, each inner operation is
identity, and the two source-predicate toggles in (A.12) cancel too.
Thus the completed $P_u$ is identity there. These statements do not
require raw dirty $T_Z$ to preserve paired sources.

One must not repair an internal truncated $T_Z$ to its full-$h$
version: with $h=2,u=e=1$, $d=(0,0)$ and $Z=(0,1)$, the prefix is
good but full $T_Z$ for $(0\ 1)$ changes the second digit and truncated
$T_Z$ does not. This toy example embeds for any actual $u<h$ by
putting the dirty selector immediately after the prefix and
using active digits throughout the prefix.
That off-prefix action is intended to cancel in the
outer word. Nor is the prefix bad set invariant after an inner bit
shear: the first shear of $(1\ 2)$ sends idle $(X,Y)=(1,1)$ to active
$(0,1)$. Its tile repair uses a different invariant, as specified in
the native appendix.
The separate top-level digit-group placement can reach ranks beyond
$e$; it retains its full-$h$ root-level charge and is not justified
by this width-$e$ prefix argument.
In the uniform $[4]^h$
address space the prefix active count is
binomial with mean $3u/4$. For $u\ge2e$ and $u<h$, a Chernoff bound gives
\[
 \Pr\{B_u<e\}\le e^{-u/24}.                            \tag{A.15}
\]
For $u=h$, eligibility $|S|\ge e$ leaves no exceptional eligible mask.
Both $T$ and the completed logical $P_u$ preserve $S$, so the
exceptional set
$E_{\rm inherited}(S)\wedge[|S|\ge e]\wedge[B_u<e]$ is invariant.
Its whole-universe fraction is bounded by the unconditional binomial
estimate; no conditional distribution given eligibility is assumed.
At that completed boundary, sorting the exceptional records by
$T(P_u^{-1}(y))$ and refilling their slots repairs it exactly;
the maps agree everywhere else. No sector is dropped. The full role recursion always
stores the complete $[4]^h$ address universe, and its row split depends
only on the independent role field, never on $S$. Thus every child
retains the same mask distribution and off-eligibility maps are identity.
Consequently, on every sector $S$, the role program applies
$C^{\otimes |S|}$ to its active digits and identity to its idle digits;
the direct sum over all masks is exactly $A^{\otimes h}$.

The prefix correction above belongs to the completed logical routing
permutation, including $T$ and $T^{-1}$ below, after all its source
conditioning and before a linear combination forms a phase. An
internal $T_{J\wedge Q}$ is not separately corrected when its surrounding
source predicate was prefix-truncated too. It never sorts the output
of the non-permutation matrix $C$.

At selected width one, let $T$ cyclically permute that eligible active
trit and be identity on ineligible masks. The mask-only transpositions
above implement $T$ and $T^{-1}$. Three temporary array copies and a
pointwise linear combination compute
$I+\eta^2(T+T^{-1})$, which is $C$ on the selected trit. Off eligibility
the two identical terms cancel, leaving the identity. The copies,
returns and cleanup add $O(V)$ at this base call. Norm two uses its
inverse via the fixed coefficient wrapper described below.

The $M=49^3$-slot basis changes in the finite frame program decompose
into a fixed number of row-addition matching batches on the first $e$
selected ranks. Nondegenerate residual spaces over $\mathbb F_3$ have
orthogonal bases: repeatedly split a nonisotropic line; its orthogonal
complement remains nondegenerate. The sign of the directed frame
difference must also be retained.

\begin{lemma}[Signed direction and one-child implementation]
\label{lem:signed-phase-edge}
Let $v$ be one such direction in standard-dot coordinates, with
$\nu=\langle v,v\rangle\in\{1,2\}$, and let
$\varepsilon\in\{1,2\}$ specify the positive or negative directed
phase. The effective child coefficient is
\[
             \lambda=\varepsilon/\nu\quad\hbox{in }\F_3.
\]
For the block direction $\langle v\rangle\otimes\F_3^f$, use one
ordinary width-$f$ child if $\lambda=1$, and one conjugated inverse
child if $\lambda=2$. Both cases retain exactly the same child volume.
\end{lemma}
\begin{proof}
Choose $A\in\operatorname{SL}_M(\F_3)$ with $Av=e_1$.
Extend $v$ to a basis and adjust the sign of another column to obtain
determinant one; $M\ge2$. A fixed elementary row-addition word realizes
$A$. Define the address pushforward by
$(Q_Af)(y)=f(A^{-1}y)$. Its translation conjugacy is
$Q_A^{-1}T_wQ_A=T_{A^{-1}w}$.
Chronologically place the addresses by $A\otimes I_f$, apply the
first-block child $(C^{\otimes f})^\lambda$, and undo the placement.
The resulting operator is
\[
 Q_{A\otimes I_f}^{-1}(C^{\otimes f})^\lambda Q_{A\otimes I_f}.
\]
Its Fourier multiplier is
\[
 \eta^{\lambda\sum_{\ell=1}^f\langle v,z^{(\ell)}\rangle^2}.
\]
The desired directed line phase has exponent
\[
 \varepsilon\sum_\ell\langle v,z^{(\ell)}\rangle^2/\nu,
\]
so the stated $\lambda$ is exact. The placement need not preserve the
quadratic form: $\nu$ is computed before placement, in the original
form. Replacing $v$ by $-v$ leaves these squares unchanged.

For $\lambda=2$, wrap the \emph{same forward child} with the
coefficientwise $\F_4$ Frobenius on both existing codec planes,
as in (A.17). This gives $(C^{\otimes f})^{-1}$ on arbitrary codec
values. The wrapper is $\F_2$-linear; its completed conjugate is
linear over the codec field. It adds two linear child-stream scans,
and no second recursive call or new coordinate. If eligibility is
false, the child is identity and the two involutions cancel there.
\end{proof}

Explicitly, the forward/inverse cases are
\[
\begin{array}{c|c|c|c}
 \varepsilon&\nu&\lambda&\hbox{child}\\ \hline
 1&1&1&\hbox{forward}\\
 1&2&2&\hbox{inverse}\\
 2&1&2&\hbox{inverse}\\
 2&2&1&\hbox{forward}
\end{array}
\]
An upward nested-frame edge has $\varepsilon=1$; a downward edge,
including $t_3\to t_4$, has $\varepsilon=2$ on its residual subspace.
The exceptional source difference splits as a positive residual
projection plus $2P_{U_a}$, so its extra unit-line child in (A.9)
also has $\lambda=2$. This flag selects an operator, not an extra
child; the rank sum $s$ is unchanged.

At selected width $e=Mf$, each rank-$a$ edge uses exactly $a$ width-$f$
children, one for each orthogonal basis direction. With $s$ the total
child count (A.9), and $W$ roles, normalized time obeys
\[
 F(e;h)\le (s/W)F(e/M;h)
       +O\!\left(\operatorname{polylog}(L+2)
          (e^\tau+\operatorname{polylog}(L+2))\right).       \tag{A.16}
\]
The additive term includes dirty-selector shears, prefix exception
repairs, basis placements and current-parent tape parking. Put
$\sigma=\log_M(s/W)$. Exact unrolling gives
\[
 O(e^{\max\{\sigma,\tau\}}\operatorname{polylog}(L+2)).
\]
Equality $\sigma=\tau$ adds only $\log e=O(\log L)$. Here the
reconstructed interchange permits $\tau=1-2^{-36}$ and $\sigma<\tau$,
so the ratio $(s/W)/M^\tau$ is strictly below one and the bound is
\[
 O(e^\tau\operatorname{polylog}(L+2)).
\]
No additional exponent
slack is required at this step.
Arbitrary active counts are
partitioned into base-$M$ power-width groups; digit-conditioned groups
partition each mask's active ranks, and each top-level placement is
performed once per group, not again in every child. Hence the complete
$A^{\otimes h}$ tensor has a strict exponent in $h$, relative to the
stated interchange and weighted-mask suppliers.

\subsection{Codec field and explicit positive word for $H_2\otimes H_2$}

The address field $\mathbb F_3$ above is not the phase field. For the
constant word use
\[
 E=\mathbb F_2[t]/(t^4+t+1),\qquad \eta=t^5=t^2+t.
\]
The quartic has no root in $\mathbb F_2$ and remainder $1$ modulo the
only irreducible quadratic $t^2+t+1$, so it is irreducible. Thus
$E=\mathbb F_{16}$, it contains $\mathbb F_4=\mathbb F_2(\eta)$, and
$\eta^2+\eta+1=0$. The four-state phase letter is
\[
 A=(\eta I_3+\eta^2J_3)\oplus[1]=C\oplus[1].
\]
It acts separately on the two $\mathbb F_4$ coefficient planes in
$E=\mathbb F_4\oplus t\mathbb F_4$. Its inverse is $A^2$.
The inverse-phase wrapper on these planes is
\[
 \widehat\sigma(a+tb)=a^2+t b^2\quad(a,b\in\mathbb F_4),       \tag{A.17}
\]
which applies the $\mathbb F_4$ Frobenius to each plane and fixes the
codec basis element $t$. It is $\mathbb F_2$-linear and is not ordinary
$\mathbb F_{16}$ squaring.
In the polynomial basis $(1,t,t^2,t^3)$ it is the explicit map
\[
 (c_0,c_1,c_2,c_3)\longmapsto
 (c_0+c_2+c_3,\ c_1+c_3,\ c_2,\ c_3).
\]
Indeed $a_0=c_0$, $a_1=c_2+c_3$, $b_0=c_1+c_2+c_3$ and
$b_1=c_3$ in $a=a_0+a_1\eta$, $b=b_0+b_1\eta$.
Three in-place XORs implement it, and repeating them gives the identity.
For $\lambda\in\F_4$,
$\widehat\sigma(\lambda z)=\lambda^2\widehat\sigma(z)$.
Consequently coefficientwise conjugation gives
\[
 \widehat\sigma C\widehat\sigma
   =\eta^2I_3+\eta J_3=C^2=C^{-1},
 \qquad
 \widehat\sigma(C^{\otimes f})\widehat\sigma
   =(C^{-1})^{\otimes f}.
\]
This holds on arbitrary $E$-valued inputs. The two passes cancel on
idle or ineligible coordinates; either pass alone can change their
coefficient values. Both passes act within the existing four-bit lanes,
in $O(V)$ time around one child. Ordinary $E$ squaring used on both
sides would fail even on an idle value $t$, sending it to
$t^4=t+1$ instead of $t$. The semilinearity assertion is only for
$\F_4$ scalars, not arbitrary $E$ scalars.

Here is the explicit word construction. Put $a=\eta^2$ and $b=\eta$;
use column vectors and rightmost-first composition. Then
\[
 A^3=I,\quad A^{-1}=A^2,\quad
 u=(1,a,a,0)^T,\quad v=(1,b,b,0)^T,\quad
 P=AE_{11}A^{-1}=uv^T,
\]
with $v^Tu=1$, hence $P^2=P$. Let
\[
 D=\operatorname{diag}(t,1,t+1,1),\quad Q=DPD^{-1}=x_0y_0^T,
\]
where $x_0=(t,a,a(t+1),0)^T$ and
$y_0=(t^{-1},b,b/(t+1),0)^T$. Direct products give
\[
 v^Tx_0=0,\qquad y_0^Tu=\kappa
 =\frac{t^2+t+1}{t(t+1)}\ne0,
 \qquad PQ=0,\quad QP=\kappa x_0v^T.                 \tag{A.18}
\]
For idempotents $P,Q$ with $PQ=0$,
\[
 [I+cP,I+dQ]=I+\frac{cd}{1+c}QP.                    \tag{A.19}
\]
This follows by using $(I+cP)^{-1}=I+c(1+c)^{-1}P$ and
$QPQ=0$. With $c=d=\eta^2$, $cd/(1+c)=1$. Thus
\[
\begin{aligned}
 G&=A\operatorname{diag}(\eta,1,1,1)A^{-1}=I+\eta^2P,\\
 H&=DGD^{-1}=I+\eta^2Q,\\
 T_0&=GHG^{-1}H^{-1}=I+\kappa x_0v^T.
\end{aligned}\tag{A.20}
\]
Each of $G,H,G^{-1},H^{-1}$ uses three forward $A$ letters after
replacing inverse $A$ by $A^2$, so $T_0$ uses 12.

The remaining conjugations isolate an elementary matrix. Set
\[
 B=\operatorname{diag}(1,\eta,\eta,1),\quad
 x=(t,1,t+1,0)^T,\quad w=(1,1,1,0)^T.
\]
Then $T_1=BT_0B^{-1}=I+\kappa xw^T$ (12 $A$ letters). Let $P_{12}$
swap coordinates 1 and 2. Since $w$ is fixed and $w^TP_{12}x=0$,
\[
 T_2=T_1P_{12}T_1P_{12}^{-1}
   =I+\kappa(t+1)(e_1+e_2)w^T                       \tag{A.21}
\]
uses 24 $A$ letters. With $D_3=\operatorname{diag}(1,1,\eta,1)$,
\[
 T_3=T_2D_3T_2D_3^{-1}
   =I+\kappa(t+1)\eta(e_1+e_2)e_3^T                \tag{A.22}
\]
uses 48; the two cross pairings vanish and the row sum is
$(1+\eta^{-1})e_3^T=\eta e_3^T$. With
$D_1=\operatorname{diag}(\eta,1,1,1)$,
\[
 T_4=T_3D_1T_3D_1^{-1}
   =I+\kappa(t+1)E_{13}                             \tag{A.23}
\]
uses 96; the cross pairings vanish and the column factor is
$1+\eta=\eta^2$. Put $\mu=\kappa(t+1)\ne0$; then
\[
 T_5=\operatorname{diag}(\mu^{-1},1,1,1)
       T_4\operatorname{diag}(\mu,1,1,1)=I+E_{13}.  \tag{A.24}
\]
All conjugating factors are invertible monomial matrices. Coordinate
permutations produce every $U_{ij}=I+E_{ij}$ with the same 96 forward
$A$ letters. This is an explicit formula, not an uncomputed word search.

Finally write $H_2=US$, with $U=I+E_{12}$ and $S$ the two-coordinate
swap. In ordered coordinates $(00,01,10,11)$, put
\[
 P_{\rm first}=(1\ 3)(2\ 4),\qquad
 P_{\rm second}=(1\ 2)(3\ 4).
\]
Then, using the four explicit elementary-matrix words,
\[
 H_2\otimes H_2
 =U_{13}U_{24}P_{\rm first}U_{12}U_{34}P_{\rm second}.       \tag{A.25}
\]
There are four $U_{ij}$ factors, hence $4\cdot96=384$ forward $A$
letters; the remaining factors expand to 392 invertible monomial
letters, for 776 total. The completed matrix has entries in $\mathbb F_2$.
As an $\mathbb F_2$-linear map on the four binary coordinates of an
$E$-symbol, it applies $H_2\otimes H_2$ independently to each coordinate
plane. Thus a binary $0/1$ embedding returns to $0/1$, and any zero lane
used for padding returns to zero at the completed word. An arbitrary
$E$-valued input need not return to the scalar subfield. Intermediate
$E$ lanes must not be cropped.
When four arbitrary payload bits are packed into one $E$ symbol,
all four coefficient planes remain data; only extra zero bits from
the final partial group are removed. No projection $E\to\F_2$
is applied to such a packed symbol.

Pairing $2h$ identical native endpoint axes and tensoring this fixed word
reduces their algebraic action to $A^{\otimes h}$ letters and fixed
monomial letters. The permutation part of each monomial is in
$\mathrm{AGL}(2,\mathbb F_2)=S_4$: its NOT/CNOT address shears and the
product of its $h$ diagonal $E$ scalars are charged by the separately
supplied native mask/routing and record-scan interfaces. The active-sector
role construction reduces the $A^{\otimes h}$ letters. The native
large-even-radix endpoint appendix supplies their tape price. No binary
$2^k$-axis H2 bound is asserted here; that requires a distinct binary
mask-routing argument and is not needed for the manuscript's native
endpoint interface.
For precision, write a monomial as $M_{\pi(d),d}=m_d$, with
$m_d\in E^*$. Apply the coefficient $\prod_i m_{d_i}$ at its current
\emph{source} address, then route by $\pi^{\otimes h}$. Equivalently
an output-address scan would use $\prod_i m_{\pi^{-1}(e_i)}$.
Using the source-indexed list unchanged after routing is generally
incorrect. Tensoring the fixed word is letterwise:
$(L_1\cdots L_{776})^{\otimes h}
 =L_1^{\otimes h}\cdots L_{776}^{\otimes h}$.
Thus there are 384 tensor phase calls, not 384 separate calls per axis.

\paragraph{Scope.}
The mask shears used here are constructed independently in
Section~\ref{sec:weighted-native-mask}; interchange is proved in
Appendix~\ref{app:swap}. The mask lemma never calls the later native
H2 tensor theorem. Conversely the rational interchange proof uses
only the scalar incidence motif and finite subspace inclusions, not
this phase recursion. These cross-references follow an acyclic
dependency order, despite the order in which appendices are printed.

Equations (A.5)--(A.10) give the finite scratch-safe routing and strict
rank saving; (A.11)--(A.16) give the active/idle and dirty-selector
interface with inherited mask eligibility; (A.17)--(A.25) give the
explicit positive word. Their tape implementation is relative to the
separately proved weighted-native-mask and equal-chunk interchange
suppliers. The appendix does not reprove those tape suppliers, the ordinary
polynomial multiplication composition, or external acceptance of the
claimed asymptotic result.


\section{Canonical endpoint kernels on mixed-radix tape addresses}
\label{app:native-endpoint}

This appendix supplies a strict-saving implementation for the canonical
adjacent-pair kernel on a product of even native radices.  It does not
implement arbitrary pair matchings, a cyclic Fourier index permutation, or
ordinary multiplication by itself.  The finite phase-word supplier and the
equal-chunk interchange contract are stated explicitly below; their
component proofs are cited separately.

\subsection{Model, assumptions, and supplied finite phase word}

Let the address set be
\[
  \Omega=\prod_{i=1}^t[n_i],\qquad n_i\text{ even},\qquad
  M=\prod_i n_i,\qquad L=\lceil\log_2(2M)\rceil,
\]
with
\[
  L^3\le n_i\le R,\qquad \log R=O((\log L)^c)
\]
for a fixed constant $c$.  Each record has $D\ge1$ arbitrary payload
bits.  The machine has a fixed finite alphabet, a fixed finite number of
one-dimensional tapes, and ordinary binary shape descriptors.  All address
work, parking, padding, and record motion are charged.  Write $V$ for the
complete current stored bit volume, including padded address fields and
opaque record lanes; on the original unpadded stream $V=MD$.

We use the interchange construction in Appendix~\ref{app:swap}, based
on the finite-network method of Section 4 of the September 2026 source
\emph{Integer multiplication below $n\log n$}: two $u$-bit
address fields in
\[
 [P]\times[2^u]\times[G]\times[2^u]\times[B]
\]
can be interchanged in $O(Vu^\tau)$ bit steps with
$\tau=1-2^{-36}<1$, where $V$ is the complete stored bit volume. This
choice follows from our exact $48$-triple reconstruction, not from
attributing that numerical exponent to the source preprint. The
prefix, gap, and suffix are arbitrary, including opaque payload.  This
appendix uses only this interchange procedure, not the paper's global
integer multiplication theorem.\footnote{Primary statement: \href{https://raw.githubusercontent.com/openai/math/4eaee87548250b2682affda7d05d947c05554bae/preprints/Integer-multiplication-below-n-log-n-September-23-2026/build/sections/04-swap.tex}{Section 4, equal-chunk interchange}. The finite role and tape reconstruction dependencies are listed at the end of this appendix.}

The second finite supplier is the exact word
\[
 A=(C\oplus[1]),\qquad C=\eta I_3+\eta^2J_3,\qquad
 \eta^2+\eta+1=0\text{ in }\mathbb F_4.
\]
The certificate proves that a fixed word in forward $A$ and invertible
monomial $4$-by-$4$ matrices over $\mathbb F_{16}$ equals
$H_2\otimes H_2$, where
\[
 H_2=\begin{pmatrix}1&1\\1&0\end{pmatrix}\quad\text{over }\mathbb F_2.
\]
It has 384 forward $A$ letters and 392 monomial letters.  This is the
finite phase-$H$ supplier used here; no generic or unproved $H_2$ tensor
theorem is assumed.  Appendix~\ref{app:finite-phase} supplies the finite
word and role proof; the earlier detailed word certificate is in
\path{supporting/ntt_base_five/H2_FINAL_FINITE_WORD_CERTIFICATE.md}.
The priced recursive tensor supplier for $A=C\oplus[1]$ is the finite
ternary role/frame construction of
\path{supporting/H2_FINAL_TERNARY_ROLE_CERTIFICATE.md}, with its dirty-router
interface detailed in \path{supporting/H2_TENSOR_CLOSED_PROOF.md}.  The role
construction is used only through the finite constants, frame identities,
and recurrence restated below.  The mixed-radix endpoint masks in this
appendix replace its binary mask router.

\subsection{Weighted two-bank tile mask}
\label{sec:weighted-native-mask}

For an integer pair $(x,w)$, define the four-update word $P(a,b,s)$ by
the chronological operations
\[
\begin{aligned}
 x&\mathrel{+}= [w\equiv b\pmod2], &
 w&\mathrel{+}=s[x\equiv a\pmod2],\\
 x&\mathrel{-}= [w\equiv b\pmod2], &
 w&\mathrel{-}=s[x\equiv a\pmod2].
\end{aligned}
\]
The parity tests are evaluated at that step.  Set
$p=P(1,1,+1)$ and $q=P(1,0,-1)$.  A direct check on the four parity
tiles gives $p=(B\ C\ D)$ and $q=(A\ B\ D)$, with
$A=(0,0),B=(1,0),C=(0,1),D=(1,1)$.  Thus the chronological word
$p^{-1}$ followed by $q$ is
\[
 (x,w)\longmapsto
 \bigl(2\lfloor x/2\rfloor+(1-(x\bmod2)),w\bigr).
\]
For a frozen control bit $g$, mask every one of these eight elementary
updates by $g$.  The resulting word is identity for $g=0$ and the displayed
bit flip for $g=1$.  In particular it restores the whole integer $w$,
not merely its parity.  This is the binary tile identity in
\path{supporting/ntt_base_five/binary_packed_affine_monomial_routing.md}.

Now take $u$ digit pairs with even radices $n_i,m_i$, all at least $L^3$,
and ordinary numeric orders
\[
 X=\sum_{i=0}^{u-1}x_iA_i,\quad A_i=\prod_{j<i}n_j;
 \qquad
 W=\sum_{i=0}^{u-1}w_iB_i,\quad B_i=\prod_{j<i}m_j.
\]
Each $g_i$ is computed solely from frozen address controls disjoint from
$X,W$ and the dock below; it cannot depend on a bank while that bank is
being updated.  In one $X$-update of the eight-step word, let
$\delta_i\in\{-1,0,1\}$ be its local signed increment, determined only
by $w_i\bmod2$, $g_i$, and the public word position.  Execute the bank
update by the single weighted integer increment
\[
                 \Delta_X=\sum_i A_i\delta_i.
\]
For a $W$-update, define $\epsilon_i$ from $x_i\bmod2$ and set
\[
                 \Delta_W=\sum_i B_i\epsilon_i.
\]
The first offset is independent of $X$ on each fixed $(W,\text{controls})$
fiber; the second is independent of $W$ on each fixed
$(X,\text{controls})$ fiber.  Hence each update is a cyclic line rotation
of its target numeric bank, not a gather of its component digits.

To implement it in the fixed-tape model, reserve a separate binary dock
$Z$ of width $c_0=\max(a,b)$, where
$a=\lceil\log_2\prod_i n_i\rceil$ and
$b=\lceil\log_2\prod_i m_i\rceil$.  Pad each target range once to
$2^a$ or $2^b$ (factor below two), retaining its true validity predicate.
For an $X$-update, interchange its $a$-bit field with an $a$-bit slice of
$Z$ using the supplied equal-chunk theorem.  On the valid prefix
$0\le X<\prod n_i$, rotate modulo the true range by $\Delta_X$; copy the
invalid tail pointwise.  If $W$ is invalid, use offset zero on the valid
$X$ prefix.  Interchange back.  The $W$-update is symmetric, forcing zero
offset when $X$ is invalid.  Every exchange restores arbitrary dirty dock
contents.  Consequently every privately padded state is fixed pointwise
by each elementary update and may be cropped only after the completed
word.  No unequal-side transpose across an arbitrary gap is presumed.
Validity follows the \emph{logical target}: after docking $X$, test
its value in the dock, not the old physical $X$ position, which now
contains arbitrary dock data. The other bank remains in place.
The internal numerical pads of each SWAP call have already been
cropped when that call returns, before this valid-prefix rotation.

Put $N_X=\prod n_i$, $N_W=\prod m_i$. Suppressing Cartesian
spectators, the original native domain is
$\Omega_0=[N_X]\times[N_W]$ and the privately padded domain is
$\widetilde\Omega=[2^a]\times[2^b]$.
Every complete docked update fixes
$\mathcal I=\widetilde\Omega\setminus\Omega_0$ pointwise.
For the concrete schedule below keep these private pads through
the repair, then crop them after the exact mask is complete.

On $\Omega_0$, define the good set $G_0$ by quotient margins
\[
 10\le\lfloor x_i/2\rfloor\le n_i/2-11,\qquad
 10\le\lfloor w_i/2\rfloor\le m_i/2-11\quad\text{for all }i.
\]
Each local integer moves by at most four during the eight updates, so on
this set no digit overflows and the weighted updates equal the independent
tile words.  The ideal completed map toggles only $x_i$'s endpoint bit;
it preserves every quotient label and fixes $W$.  It therefore preserves
$G_0$ and $\mathcal B=\Omega_0\setminus G_0$ separately.
Extend the ideal map as identity on $\mathcal I$. The packed word
$P$ agrees with that extension on $G_0\cup\mathcal I$, and both
are permutations of $\widetilde\Omega$. Hence
$P(\mathcal B)=T(\mathcal B)=\mathcal B$.
The union bound gives density on the \emph{original native domain}
\[
 \delta=\frac{|\mathcal B|}{|\Omega_0|}
 \le40\sum_i(n_i^{-1}+m_i^{-1})=O(u/L^3).
\]
The repair marker is explicitly
\[
 [X<N_X]\wedge[W<N_W]\wedge[\text{a quotient margin fails}].
\]
It is false on $\mathcal I$, even if private padding occupies a
constant fraction of $\widetilde\Omega$. Do not decode an invalid
numeric bank as a wrapped native digit tuple.
Extract only marked records at their current padded output addresses $y$. Recover
their sources $x=P^{-1}(y)$ by reversing the eight packed updates, then
compute the ideal endpoint destination $T(x)$. These are public address
calculations with frozen controls; no old key is assumed to accompany
the payload. Keys are destination ordinals in the same current padded
layout. Stably sort these destination keys and refill the same bad
slots.  This corrects every exceptional address and discards no payload.
The repair costs $O(\delta VL+M_{\rm rec}L^C)$ for a fixed $C$, where
$M_{\rm rec}$ is the native record count before these private bank
pads. The padded record count and bit volume are each below four
times their native values, so full marking/refill scans cost $O(V)$.
For clarity, with $K\le\delta M_{\rm rec}$ bad records, use stable
binary radix sorting on the $b=O(L)$ key bits. In each pass copy one
fixed-width key/payload item to a staging tape, remember the selected
bit, rewind, and copy to the appropriate zero/one bucket. Concatenate
the buckets and rewind for the next pass. A fixed number of tapes
and $O(K(R_{\rm op}+b))$ steps suffice per pass, including all item
and file returns. Counter updates are amortized over the scanned
length, and resetting a counter costs at most the following item
or block length. The bound is therefore
$O(KL(R_{\rm op}+L))$, plus explicitly charged polynomial address
work and full-array marking/refill scans. At refill, consume the next
sorted item exactly at a marked bad slot;
$P(\mathcal B)=T(\mathcal B)=\mathcal B$ proves the
order without a full-key comparison at every good record. The
validity-and-margin predicate above marks $\mathcal B$ directly, so no extra mask lane is added
to the retained field codec. The exponential opaque width
below absorbs the key-only and metadata terms. The tile word uses
a constant number of dock interchanges and rotations, hence costs
$O(V(u\log R)^\tau+V)$.  In the application below, the bundled opaque
record width is exponential in $L$, so it absorbs $M_{\rm rec}L^C$; also
$\delta L=O(u/L^2)=O(1)$.  The resulting mask and its inverse cost
\[
                       O(V(u\log R)^\tau+V).
\]
After refill the padded map is exactly $T$ on $\Omega_0$ and identity
on $\mathcal I$, with the dock restored. Crop only the private $X,W$
numeric tails now, returning the original native ranges before any
enclosing phase child. A caller's Fourier dummy $p_i$ in a native
range $n_i=p_i+1$ is \emph{inside} that range and is not such a tail;
it remains part of the native kernel. Role-row, suffix and coefficient
bit pads also retain their separate outer crop boundaries below.

\subsection{Native pair endpoints and disjoint banks}

For a native digit $x_i\in[n_i]$, write $x_i=2a_i+b_i$ with
$b_i\in\{0,1\}$.  The desired local matrix applies $H_2$ to
$\{2a_i,2a_i+1\}$ for each fixed $a_i$.  Equivalently, for every fixed
tuple $(a_1,\ldots,a_t)$ it applies $H_2^{\otimes t}$ to the endpoint
tuple $(b_1,\ldots,b_t)$, while the physical interleaved mixed-radix
record order remains unchanged.  This is a statement about fibers of the
native address set; it does not perform or price a tensor regrouping or an
endpoint gather.  The labels $a_i$ are spectators.  Because every $n_i$
is even, endpoint bits are independent uniform bits under uniform native
addresses, and pairs of endpoint bits from disjoint axes are uniform
four-state digits.  No arbitrary matching is hidden in this factorization.

For a large block of $2h$ endpoint axes, use the already contiguous
product rank of native intervals; no per-record mixed-radix conversion is
performed.  One aggregate rectangle exchange moves the chosen interval
to the front.  For adjacent ranges $[A][B]$, pad once to
$A'<2A,B'<2B$, perform the binary rectangle transpose, and crop the
completed rectangle before returning.  Its volume is below four times
the original and its charged cost is $O(VL^\gamma)$, where
\[
                  \gamma=\frac1{2-\tau}<1.
\]
This follows from Euclidean decomposition into equal-chunk interchanges;
partial internal exchanges need not preserve the valid rectangle.  Reverse
the complete transpose after the block.  This is a top-level layout cost,
not a cost paid in each recursive role child.

Choose
\[
 \ell_R=\lceil\log_2 R\rceil,\qquad
 h=\min\!\left(\left\lfloor t/16\right\rfloor,
                 \left\lfloor L/(40\ell_R)\right\rfloor\right).
\]
For sufficiently large $L$, $h\ge1$ and $5h<t/2$.  The first two
disjoint $h$-axis banks are $X,Y$.  From the remaining consecutive native
interval take three more disjoint $h$-axis banks $Q,Z_0,W$; $Z_0$ here is
an endpoint bit bank, distinct from the private binary dock $Z$.  The
other native axes form a spectator suffix.  The F16 word works on
four-state digits $d_j=2b_j+b_{h+j}$, pairing the $j$th axis of $X$
with the $j$th axis of $Y$; $Q,Z_0,W$
are dirty endpoint-selector banks.  This is logical bank indexing only:
the bits remain in their two contiguous native banks, and each shear is
executed by the mask oracle without physically interleaving $X$ and $Y$.
All controls for one weighted mask
are chosen from the other disjoint banks and role bits, never its target,
tile auxiliary, or dock.

More explicitly, the dirty-router shears for
$(0\ 1):Y\mathrel{\oplus}=Z_0\wedge\neg X$ use $X,Z_0$ as frozen
controls.  The three shears for
$(1\ 2):X\mathrel{\oplus}=Z_0\wedge Y$,
$Y\mathrel{\oplus}=Z_0\wedge X$,
$X\mathrel{\oplus}=Z_0\wedge Y$ each use the other plane and $Z_0$.
The outer source-value selector $F_a$ for $Q$ is computed from $X,Y$;
the inner selector for $Z_0$ is $J(S)$ or $J(S)\wedge Q$, computed
from $X,Y,Q$. Each is reversed after use at its specified completed
boundary.
The tile auxiliary for every call is a slice of the separate native $W$
bank, disjoint from the current target and every control.  Role bits may
be added as frozen controls.  These are the control dependencies of the
dirty-selector construction in
\path{supporting/H2_TENSOR_CLOSED_PROOF.md}, Lemma 2; no selector is recomputed
from a bank while that bank is the weighted tile auxiliary.

Pad the numeric range of the whole native spectator suffix once to its
next power of two, a factor below two.  Split this already contiguous
binary index into high role bits, then dock bits, then low bits.  Write
$r$ for the number of role bits, $s$ for the padded suffix width, and
\[
 c_0=\max_{G\in\{X,Y,Q,Z_0,W\}}\lceil\log_2|G|\rceil
       \le h\log R+1
\]
for the reserved dock width. The high row field is borrowed from these
existing suffix index bits, not added as another Cartesian factor.
Reserve there at least
$W_*^{\lceil\log_m h\rceil}$ rows, padded once by a factor below two as
in the finite role supplier; each recursive level divides this field by
$W_*$.  Follow it by the $c_0$ dock bits.  The remaining low suffix index is bundled with the
external payload.  Since $5h\log R\le L/8$, the suffix has width at
least $7L/8-O(1)$; $r=O(\log h)$ and $c_0\le L/40+1$.  Thus for all
sufficiently large $L$ the opaque record width obeys
\[
 R_{\rm op}=2^{s-r-c_0}D\ge 2^{4L/5}D.
\]
Move only the $r$ high role bits to the front by individual binary
digit moves, costing $O(Vr)=O(V\log h)$, and undo these moves once
the complete block returns. The dock is now the rightmost active
address field, immediately before the low opaque lanes. For a maximum
role depth $k$, choose $r$ with $2^r\ge W_*^k$, and pad this one
independent row range to $W_*^k\lceil2^r/W_*^k\rceil<2^{r+1}$.
The row pad stays outside the role recursion. In the program used here
it is retained through the entire fixed H block, including all phase
and monomial letters, and is cropped after that completed word restores
row labels. Monomial letters act only on $X,Y$ and codec payload, so
they preserve this independent numeric row coordinate.
Here $r$ counts borrowed binary suffix bits \emph{before} this private
padding. After padding, the row range is a separate numeric factor
$R'$, which need not fit in $r$ bits. Its counter can use $r+1$ bits;
this is metadata, not another bit borrowed from the low opaque index.
The array volume grows by $R'/2^r<2$, already included in the current
$V$. Role splitting uses numeric rows and SWAP permits an arbitrary
numeric prefix, so neither requires an $r$-bit encoding of $R'$.
Crop the row range back to its original $2^r$ values before undoing
the $r$ binary digit moves. Thus no genuine suffix states are dropped,
and the exponent $s-r-c_0$ for the original low opaque lanes is unchanged.
All polynomial-in-$L$ selector generation and repair metadata is paid by
this width, even when $D=1$.  No control reads the low suffix index or
the original payload.

\subsection{A fixed metadata program and sufficient suffix threshold}

Here is a concrete bound for the instruction family actually used,
rather than an allowance for every possible polynomial-time function.
Decode the five bank indices by at most $5L$ schoolbook divisions on
$L$-bit integers, at cost $O(L^3)$. Store the active-position list,
pair matching and any group-placement permutation in at most $L$
records of $O(L)$ bits each. Charge $O(L^2)$ for a full-scan lookup;
there is no unit-cost indexed access. A fixed-slot basis instruction
joins ranks $(i-1)f+a$ and $(j-1)f+a$, for $a=1,\ldots,f$, with
distinct fixed slots $i,j$. A group placement is an interval rotation
of the active list; its cycle decomposition uses at most $L$ seeds,
$L$ steps per seed and full-scan lookup/marking. The two cycle
reflections then give its disjoint swap lists.

The recursive descriptor retains a list of $O(\log h)$ mask-only
eligibility conditions, not an expanded recursive Boolean expression.
They are conjunctions of active-count, prefix and base-$m$ digit tests;
all completed coordinate-value maps preserve the physical active mask.
Every actual Boolean-vector instruction is one of $F_a(X,Y)$,
$J(S)$, $J(S)\wedge Q$, $Z_0\wedge Y$, $Z_0\wedge X$,
$Z_0\wedge\neg X$, or a fixed NOT/CNOT. A constant number of the
same routines computes each weighted offset and each inverse packed
update used for $T(P^{-1}(y))$. Field scalars are products of at most
$h$ entries from fixed $\F_{16}$ tables. These are fixed finite
compositions of loops with bounds at most $L$ or $L^2$. Allowing
redundant list copies and descriptor scans gives the deliberately
loose bound $O(L^{16})$ time and $O(L^4)$ metadata storage per current
address. Copying the opaque word and moving sort payloads remain
separate, explicitly linear-volume charges.

Fix $P_{\rm nat}(L)=L^{64}$. It pays this actual generator whenever
$R_{\rm op}\ge L^{64}$. For an explicit suffix choice let
$k=\lceil\log_m h\rceil$ and $r=56k$, where $m=117649>2^{16}$.
The fixed phase role count below obeys
$W_*<18000^3\cdot9000<2^{56}$, so $2^r\ge W_*^k$ and
$r\le56+(7/2)\log_2 L$. Since $s\ge7L/8-2$ and
$c_0\le L/40+1$,
\[
 s-r-c_0\ge(17/20)L-(7/2)\log_2 L-59.
\]
For $L\ge2^{16}$ this is at least $(4/5)L$, and
$2^{(4/5)L}\ge L^{64}$. Thus the original suffix supplies the fixed
metadata allowance for every external $D\ge1$. The privately padded
numeric role range still adds less than two volume, without consuming
another original suffix bit. This certificate concerns nonempty
admissible five-bank blocks; the direct finite branch also covers
smaller $L$ and parameters for which $h=0$.

One may fix the prefix constant to $4096$:
\[
 u=\min\{h,\lceil4096(e+\lceil\log_2(L+2)\rceil)\rceil\}.
\]
If $u<h$, the prefix exception fraction is at most
$\exp(-u/24)\le(L+2)^{-128}$; if $u=h$, eligible masks have no
prefix exception. More precisely the repaired mask is
$E_{\rm inherited}(S)\wedge[|S|\ge e]\wedge[B_u<e]$;
off-eligibility sectors already have the correct identity action and
are not sent through a full-payload correction sort. Its exact
permutation repair has $O(V)$ payload
cost. The separate tile-boundary fraction is at most $80u/L^3$ and
also gives $O(V)$ repair. These statements count addresses in the
retained universe and leave no probability of an incorrect output.
The two repairs have different boundaries. Tile repair follows one
completed eight-update endpoint XOR and uses invariant quotient
labels. Prefix repair follows a completed logical active-value route,
with its source conditioning and dirty selectors restored; only that
whole route agrees with its intended full-rank permutation on good
prefixes. No prefix repair is applied to an internal raw $T_Z$ or bit
CNOT. The prefix implementation is identity on the exceptional masks,
so its repair key there is simply the ideal logical destination.
Both repairs precede the linear combination defining $C$ and leave
the retained codec and opaque width unchanged.

\subsection{Finite phase recursion with inherited masks}

Group corresponding endpoint bits into $h$ four-state digits.  State $3$
is idle, and the active mask is
$S=\{j:d_j\ne3\}$.  The finite ternary role supplier implements
$A^{\otimes h}$, $A=C\oplus[1]$, by active-rank frames, dirty-selector
commutators, and a fixed finite list of matching SUM batches.  The finite
frame counts are restated here to specify the exact recurrence.  At
ground-set parameter $g=48$, put
\[
\begin{aligned}
 v&=\binom{48}{3},\quad z=3\binom{45}{2},\quad N=v^3,\quad I=3v^2,\\
 d&=49,\quad m=d^3=117649,\\
 W_*&=2N+I(vz+48),\quad L_*=I\cdot48\cdot49,\\
 s_*&=W_*m-N+2L_*.
\end{aligned}
\]
There are $W_*$ row-role streams and $s_*$ recursive children.  Every
child has volume $V/W_*$.  The source rank-one endpoint correction is
included in $s_*$.  The exact integer certificate gives
\[
\begin{aligned}
 \frac{W_*}{N}&=\frac{9633881}{1081},\qquad
 \frac{L_*}{N}=\frac{441}{1081},\\
 1-\frac{s_*}{W_*m}&=\frac{199}{9633881\cdot117649},\\
 \sigma&=\log_m(s_*/W_*)<1-2^{-36}=\tau.
\end{aligned}
\]
Indeed $1-s_*/(W_*m)=(N-2L_*)/(W_*m)$.  The finite certificate uses
$m<2^{17}$, hence $\ln m<17\ln2<12$, and
$199\cdot2^{36}>12\cdot9633881\cdot117649$ to obtain the displayed
strict exponent without a floating-point estimate.  The H=48 phase,
frame, endpoint, and source-count derivation is in
\path{supporting/ntt_base_five/h2_from_chunk_swap_integrated_audit.md}
(especially its Sections 3--4) and
\path{supporting/ntt_base_five/finite_motif_rank_closed_form_audit.md}.

\paragraph{Prepared recursive entry.}
The recursive routine is $\operatorname{PhasePower}(e,h,R_{\rm row},B_E)$,
with the bank/dock/codec descriptors and inherited mask eligibility
also passed. The maximum native radix remains $R$; $R_{\rm row}$
denotes a distinct numeric row range. The entry requires $e=m^k$,
with $R_{\rm row}$ divisible
by $W_*^k$, and the already prepared five full $h$-digit native banks.
Here $B_E=4\lceil R_{\rm op}/4\rceil$ is the existing opaque codec
word, containing arbitrary $E$ values. Its actual input volume is
\[
 V=R_{\rm row}B_5 2^{c_0}B_E,
\]
where $B_5$ is the product of the five native bank ranges. An edge
direction invokes
\[
 \operatorname{PhasePower}(e/m,h,R_{\rm row}/W_*,B_E)
\]
directly after its basis placement. The child keeps the same physical $h$, bank ranges,
dock, opaque width and codec; its row range is divisible by
$W_*^{k-1}$ and its volume is exactly $V/W_*$. This entry performs
no bank layout, suffix padding, row padding, or new field embedding.
Its return has the same prepared shape and can be nonbinary.
For a directed residual phase with sign $\varepsilon$ and direction
norm $\nu$, record $\lambda=\varepsilon/\nu\in\{1,2\}$ as in
Lemma~\ref{lem:signed-phase-edge}. Use the ordinary child for
$\lambda=1$; $\lambda=2$ wraps this one call on both existing $\F_4$
planes, with no additional child or payload-plane split. Thus downward
unit-norm directions and the source correction use the inverse, while
downward norm-two directions use the forward child.
The explicit four-bit map after (A.17) is applied pointwise to the
current child input and output, in two $O(V/W_*)$ scans. It changes
no address, row, bank, dock, codec width, or parked parent data.
Conjugacy uses the completed child's $\F_4$ phase coefficients;
the wrapper is not an automorphism for arbitrary $E$ multiplication.

Only a root preparation routine constructs the independent row and
codec. The suffix, row and conservative codec factors are respectively
below two, below two, and at most four, so its prepared volume is
at most $16MD$ after aggregate rectangle pads have been cropped.
The routine for a full phase letter, $\operatorname{PhaseAll}$,
calls the exact power routine on its base-$m$ groups. A power child
does not re-enter $\operatorname{PhaseAll}$ or root preparation.
Explicitly write the current active count as
$a=\sum_j a_jm^j$. For each fixed $j$ and $i=0,\ldots,m-2$,
use eligibility $a_j>i$, group width $e=m^j$, and group start
$\sum_{\ell<j}a_\ell m^\ell+i m^j$. Move that consecutive active
group to the first $e$ active ranks, call the prepared power entry,
and undo its placement. For each mask these groups concatenate,
without overlap, to its $a$ active ranks. All passes keep the same
numeric row and codec.

Each completed logical active-value route preserves $S$ and restores
its dirty auxiliaries. Inner bit shears need not preserve $S$; their
separate tile correction uses quotient labels instead. A fixed monomial
can change $S$, but it is itself an exact
address permutation; recompute the active mask for the next role pass.
The row-role split is independent of $S$.  Hence every child retains the
complete physical mask distribution and the same role multiplicity.
Eligibility is a mask-only predicate with the active-count test at
entry; off-eligible frames act as identity. Prefix concentration is applied to the current
mask: for
\[
 u=\min\{h,\lceil C_0(e+\log(L+2))\rceil\},
\]
the first $u$ four-state digits contain at least $e$ active ranks except
with probability at most $\exp(-u/24)$, for a sufficiently large fixed
$C_0$. The bad-prefix sector is extracted only at a completed logical
route boundary, after all nested source conditioning. At that boundary
the implemented route equals the ideal on good prefixes and is identity
on the exceptional masks; its dirty selectors have returned. The
ideal router preserves the whole active mask $S$, so this exceptional
sector is invariant and sorting/refill by the ideal destination is
exact. The weighted
tile guard exception from the preceding subsection is repaired separately;
neither correction crops or drops a record.

\paragraph{Induction on the full universe.}
For each physical mask $S$, the width-$e$ contract is
$C^{\otimes e}$ on its first $e$ active ranks if the inherited predicate
and $|S|\ge e$ hold, and identity otherwise. The child receives this
complete eligibility test, including the parent's threshold; checking
only $|S|\ge e/m$ would be insufficient. Fixed basis placements act on
these active values and preserve $S$ even when the selected ranks are
not physically contiguous. By induction each edge direction has its
specified signed phase on arbitrary $E$-valued payloads. For
$\lambda=2$ the two-plane Frobenius conjugation completes to the required
$E$-linear inverse phase. The finite common-label XOR program then gives
$C_{\rm out}C_{\rm in}^{-1}=C_{\rm full}$ on every routed role,
including arbitrary scratch roles. Its inverse role merge removes the
scalar role permutation by \eqref{eq:full-role-phase}.
Off eligibility every phase is identity and
the same merge returns identity. No sector is extracted and no clean
payload value is used in this induction.

Replace each binary frozen XOR mask in the finite router by the weighted
two-bank endpoint mask.  A selected prefix of $u$ endpoint digits has
numeric width at most $u\log R+O(u)$.  The reserved dock is borrowed only
in a slice wide enough for the current target prefix, and is returned
after each elementary update.  The whole $c_0$ dock is never exchanged
inside a small child. The one-node toll divided by its current volume
$V$, including exact prefix and tile exception repairs, is bounded by
\[
 O\!\left((\log L)^C\bigl((e\log R)^\tau+(\log L)^C\bigr)\right)
\]
for a fixed $C$.  Metadata scans cost $O(M_{\rm rec}L^C)$ and are absorbed by
$R_{\rm op}$.  Parking the other current role streams costs $O(V)$ per
child, a fixed additive node toll because $W_*,s_*$ are constants.
Descendants never scan parked ancestor streams.

To implement that last assertion, use a fixed set of role tapes, a
parking stack and a descriptor stack. Before one child, append only the
other current role streams and separators to the parking top, copy the
active stream to child I/O, and clear the current work intervals.
After the child returns, restore inactive streams by popping right to
left and writing right to left into prepared intervals. This preserves
their internal order and ends at the saved current-parent boundary.
The child's own stack material has already been popped. No operation
crosses a parked ancestor. The total passing/return cost per child is
$O(V)$ at its parent, and the number $s_*$ of children is a fixed
constant. It adds a fixed normalized toll, without multiplying child
time. All children keep the full five-bank native coordinate universe and the
same low opaque lanes, so polynomial-in-$L$ descriptor work fits in
their own volume even when their row count is small. SWAP uses a
separate fixed tape family and its own stack; it never calls this
phase recursion. The total number of tape families stays finite.

For exact power widths $e=m^k$, with bounded $k$ handled by direct finite
tables, the normalized recurrence is therefore
\[
 F(m^k)\le \frac{s_*}{W_*}F(m^{k-1})
       +O\!\left((\log L)^C
          \bigl((e\log R)^\tau+(\log L)^C\bigr)\right).
\]
For an arbitrary active count, the supplier partitions the ranks into
base-$m$ digit-conditioned consecutive power-width groups and handles
those groups in separate passes; their widths form a geometric sum.
The group placements are charged separately at full block width.  There
is no physical gather implicit in a group placement. For its public
active-coordinate permutation, decompose each cycle $j\mapsto j+1$
as the two reflections $j\mapsto-j$ and $j\mapsto1-j$. Each reflection
is a disjoint matching of coordinate swaps computable from the whole
mask $S$. On one pair, the chronological trit updates
$y\gets y+x$, $x\gets x-y$, $y\gets y+x$, $x\gets-x$ give the swap.
Each shear is a disjoint source/target matching from the selector
construction of Appendix~\ref{app:finite-phase}, and all masks and
quotient labels return at its completed boundary. Thus a placement
and its inverse are priced at full $h$ width, in
$O(V(h\log R)^\tau\operatorname{polylog}L)$ time. There are only
$O(\log h)$ digit-conditioned passes, with a fixed $m$-dependent
constant. These are root placement tolls, rather than costs incurred
at each smaller child. There is no ceiling recurrence for arbitrary
$e$. Put $a_0=s_*/W_*=m^\sigma$ and absorb the fixed logarithmic
radix factors in $Q(L)=(\log(L+2))^{C'}$. The exact power recurrence is
\[
 F_k\le a_0F_{k-1}+O(Q(L)(m^{k\tau}+1)),\qquad F_0=O(Q(L)).
\]
Since $\sigma<\tau$, its geometric ratio $a_0/m^\tau$ is strictly
below one, uniformly in $L,h$ and the current row range. Thus
$F_k=O(m^{k\tau}Q(L))$. More generally the exponent is
$\max\{\sigma,\tau\}$, with an extra factor $k=O(\log L)$ only
when they coincide. Summing the separately charged groups and their
full-$h$ placements gives $F(h)=O(h^\tau\operatorname{polylog}L)$.
This uses the
finite phase-$H$ supplier above, not an assumed generic fast $H_2$ tensor.

Apply the 776-letter word, containing 384 phase letters, to $h$
four-state digits. The permutation part of each monomial
is an affine permutation of $\mathbb F_2^2$; its tensor power is a
constant number of NOT/CNOT mask stages, implemented by the same weighted
router. If $M_{\pi(d),d}=m_d$, first scan the source addresses and
multiply by $\prod_i m_{d_i}$, then perform that routing. This costs
$O(L^{16}+h+R_{\rm op})$ per address: the already certified public
address decoder, $O(h)$ fixed-field products for its coefficient,
and a linear scan of the opaque record. Here $h\le L$
and $R_{\rm op}\ge L^{64}$ in the admissible range, so the scalar
stage costs $O(V)$. No stale source-indexed coefficient list is
applied at a post-permutation address. Pad the payload to fixed
$\mathbb F_{16}$ lanes; the complete word has binary entries and returns
the auxiliary lanes to zero.  Do not crop these lanes inside an $A$
letter, monomial, or intermediate scalar stage.

The completed block applies $H_2$ to the selected $2h$ endpoint axes and
identity to every other native label, endpoint, role, dock, and suffix
index.  Its ordinary binary transpose is the same completed operator,
since $H_2$ is symmetric.  Its inverse is two forward applications,
since $H_2^3=I$.

\subsection{Pad lifetimes, block sum, and exponent}

The distinct pads have disjoint lifetimes:
\begin{itemize}
 \item Each private $X$ or $W$ numeric pad is fixed pointwise during all
   eight updates, excluded from the exception sort, and cropped only
   after its completed weighted mask including exact repair.
 \item Each aggregate native-rectangle pad is cropped after the complete
   rectangle transpose, before entering the role recursion; the inverse
   aggregate transpose is applied after the block.
 \item The high role-row pad is retained through the whole fixed H block.
   Phase calls return its numeric coordinate, and monomial letters leave
   it unchanged. It is cropped after the binary word is complete.
 \item The low suffix pad and bundled $\mathbb F_{16}$ lanes are retained
   through the complete block word.  The block acts as identity on those
   suffix indices and has binary entries, so their artificial values
   return to zero only at that boundary.
 \item The physical dock is restored after each elementary weighted
   update.  It is not a recursion child field and never accumulates state.
\end{itemize}
No factors from these pads are reapplied in the recursive child volume.
At the outer binary-word boundary, crop the numeric row to its original
$2^r$ values and strip the codec back to $R_{\rm op}$ bits before
undoing the $r$ role-bit moves. Then crop the whole spectator suffix
and undo the bank-layout exchanges. This order never treats a padded
numeric row as an $r$-bit field, or a padded codec as original index bits.

Partition the $t$ endpoint axes into blocks of at most $2h$; a short final
block uses the same supplier at its smaller width, and one unpaired axis
is a direct two-row map.  Because $h\le t/16$, each block has the required
three disjoint selector banks. For sufficiently large $L$,
$L-2<\log_2M\le t\log_2R$ makes the second term in the definition of $h$
the smaller one. Thus
$h=\lfloor L/(40\ell_R)\rfloor\ge L/(80\ell_R)$.
Since also $t\le L/(3\log_2L)$, the number of blocks is
$O(\log R)=\polylog(L)$.  Their completed operators commute and restore
all borrowed indices.  Initial/final aggregate layouts cost
$O(VL^\gamma)$ per block.  Summing block and role costs gives
\[
 O\!\left(VL^{\max\{\sigma,\tau,\gamma\}}(\log L)^{C''}\right).
\]
With $\tau=1-2^{-36}$, $\sigma<\tau$ and
$\gamma=1/(2-\tau)>\tau$, the same construction therefore gives
\[
                         O(VL^\gamma\operatorname{polylog}L).
\]
Retaining this parameter avoids repeated numerical weakenings at later
compiler interfaces. More generally their exponent is
$\max\{\sigma,\gamma\}$. Bounded remaining parameters use direct
finite maps and enlarge the single uniform constant.

\subsection{Dependencies and exact scope}

The construction depends on:
\begin{enumerate}
 \item the interchange construction of Appendix~\ref{app:swap}, with its fixed
   tape and arbitrary-gap interface;
 \item exact stable extraction/sorting of public exceptional address
   keys, with its full-payload merge-sort bound shown above and further
   implementation details in
    \path{supporting/OPAQUE_RECORD_PACKED_ROUTING_AUDIT.md} and
    \path{supporting/H2_FINAL_PACKED_ROUTER_CERTIFICATE.md};
 \item the weighted tile identity and safe-set/repair proof above,
   with its earlier expanded ledger in
   \texttt{MIXED\_RADIX\_WEIGHTED\_TILE\_MASK\_LEMMA\_20261030.txt};
 \item the finite $A=C\oplus[1]$ role/frame supplier and the certified
   $H_2\otimes H_2$ word proved in Appendix~\ref{app:finite-phase}.
\end{enumerate}
The Euclidean aggregate-rectangle schedule and its fixed-tape return
protocol are given in
\path{supporting/ntt_base_five/fixed_tape_primitive_reconstruction.md} and
\path{supporting/ntt_base_five/chunk_swap_uniform_completion_audit.md}.
The weighted tile supplier's control bits must be frozen during its
eight updates and disjoint from target, auxiliary, and dock.  The
canonical bank schedule above provides that disjointness for every
individual shear; a selector that reads its changing auxiliary would
fall outside the lemma.  The supplied role proof must also be read with
its exact source correction and pad lifetimes; a generic circuit-size
bound does not imply this tape recurrence.

This appendix proves only the canonical adjacent-pair tensor on the
stated mixed-radix family using these explicit constructions.
It gives no cheap conversion for an arbitrary perfect
matching on a native digit.  It also does not price an ordinary cyclic
DFT-to-prime-grid permutation or the scalar arithmetic required for a
full multiplication construction.  Those are separate interfaces and
must not be inferred from the endpoint theorem.




\section{Common fixed-field words and exact-width replay}

This appendix makes the local characteristic-two substitution explicit.
It does not import the complex convolution service from \cite{local}.

\paragraph{Newton symmetry over a ring.}
Let $P$ have as column $k$ the monomial coefficients of
$\prod_{a<k}(Z-u^a)$. It is unit upper triangular, and evaluation
gives $F_p(u)P=N$. With $D=\operatorname{diag}(N_{kk})$ and
$L=ND^{-1}$, this is the unit-lower/diagonal/unit-upper factorization
$F_p(u)=LDP^{-1}$. Such a factorization with invertible diagonal is
unique over a commutative ring: comparing two makes a unit lower
triangular matrix equal an upper triangular one, identifying the
diagonal and then the two unit triangular factors. Since $F_p(u)$ is
symmetric, its transposed factorization is another of that form.
Therefore $P^{-1}=L^T$ and $F_p(u)=ND^{-1}N^T$. The pivot inverses
are exactly the units proved in the main text; no field cancellation
has been assumed.

\paragraph{A shared eight-$H$ schedule.}
Write $D_a=\operatorname{diag}(a,1)$, with $a=\eta$ or $\eta^2$.
Direct multiplication in $\F_4$ gives
\[
 H^3=I,\qquad
 H D_\eta H^2 D_{\eta^2}H
   =\operatorname{diag}(\eta^2,\eta)J,
\]
and
\[
 W_4:=H D_\eta H D_{\eta^2}H D_\eta H D_{\eta^2}=I.
\]
Thus $I=W_4W_4$, $J$ is its displayed four-$H$ word with a leading
unit diagonal followed by $W_4$, and $H=H^4W_4$ all have eight $H$
letters with unit diagonals in common slots. Identity diagonals fill
missing gaps. The formulas include all identity pairs and never
select a zero scalar. Coefficients may depend on a pair label because
every intermediate letter preserves that label. Controlled
coefficients are evaluated only at the completed $K_T$ boundaries
described in the main proof.

\paragraph{Formal convolution over a coefficient algebra.}
For comparison, Cantor--Kaltofen \cite{cantor-kaltofen} give sharper
classical size bounds and logarithmic arithmetic depth. Their sparse
integer \emph{stage} matrices compose to a formal bilinear identity;
the composed matrices need not be sparse. We use the following
elementary triadic construction instead, so its descriptor generation
and coefficient types can be checked directly.

\begin{lemma}[Elementary formal convolution]
\label{lem:formal-convolution}
Over any commutative $\F_2$-algebra $A$, ordinary multiplication of
two length-$v$ polynomials has a division-free formal bilinear circuit
with size $O(v\log^8(2v))$, depth $O(\log(2v))$, binary fan-in,
and deterministic polynomial-$v$ wire-descriptor generation.
\end{lemma}
\begin{proof}
Use schoolbook products below the fixed threshold $v_0=6^6$, with
balanced output sums. Otherwise let $b$ be the least power of three
with $b^2\ge v$, put $s=\lceil v/b\rceil$, and let $t$ be the
least power of three with $t\ge2s-1$. Then
\[
 \sqrt v\le b<3\sqrt v,\quad s\le b,\quad
 t\le3b,\quad t<6s,\quad t\mid3b.
\]
Work formally in $R=A[Y]/(Y^{2b}+Y^b+1)$, free of rank $2b$
over $A$. The polynomial is $\Phi_{3b}$ reduced modulo two.
Its roots over $\overline{\F}_2$ have exact order $3b$: they
satisfy $Y^{3b}=1$, and none satisfies $Y^b=1$ since the defining
polynomial then equals one. Every proper divisor of $3b$ divides
$b$. Therefore $Y^j-1$ is coprime to this polynomial for
$0<j<3b$. The corresponding Bezout identity extends from $\F_2$
to $A$, proving principality of $Y$ and of
$\zeta=Y^{3b/t}$ at order $t$. The inverse DFT uses $\zeta^{-1}$;
its normalization is one because $t$ is odd.

Encode the input as $\sum_{i<s}A_i(Y)Z^i$, with
$\deg A_i<b$, and similarly for the other operand. Take their
$t$-point DFTs over $R$. At a frequency the value is \emph{one}
$R$ element, represented by $2b$ coefficients from $A$; no
unevaluated $Z$ variable remains. Multiply two such values by one
recursive ordinary length-$2b$ product over $A$, and reduce the
degree-$\le4b-2$ result modulo $Y^{2b}+Y^b+1$.
For $0\le r\le2b-2$, the high monomial reductions are
\[
 Y^{2b+r}=
 \begin{cases}
   Y^{b+r}+Y^r,&r<b,\\
   Y^{r-b},&b\le r\le2b-2.
 \end{cases}
\]
They cost $O(b)$ additions in constant depth.

The inverse DFT recovers
$C_j(Y)=\sum_{i+k=j}A_i(Y)B_k(Y)$ without wrap in $j$, since
$t\ge2s-1$. Each true $C_j$ has degree at most $2b-2$, below the
quotient degree, so reduction has lost no coefficient information.
Reconstruct the ordinary product as $\sum_j C_j(X)X^{bj}$:
adjacent length-$2b$ blocks overlap, and their overlaps are added,
not concatenated. Crop the known padded zero tail to length $2v-1$.

Multiplication by a root power in $R$ pads a $2b$-coordinate
vector to $3b$, rotates, and folds
$(a_0,a_1,a_2)\mapsto(a_0+a_2,a_1+a_2)$.
Each coordinate uses at most two inputs. A radix-three butterfly
therefore has $O(b)$ binary additions and constant depth over $A$.
Explicitly, at transform length $t=3u$ split inputs as $x_{3j+r}$,
$r=0,1,2$. Compute their three length-$u$ transforms $B_r(q)$
with root $\zeta^3$, then print
\[
 \widehat x_{q+\ell u}
   =\sum_{r=0}^2\zeta^{r(q+\ell u)}B_r(q),
 \qquad 0\le q<u,\quad \ell=0,1,2.
\]
This recursively specifies every butterfly and twiddle index.
The two forward and one inverse transforms, pointwise folds and
reconstruction give, for an effective absolute constant $C$,
\[
 S(v)\le tS(2b)+Ctb\log(2t),\qquad
 D(v)\le D(2b)+C\log(2t).
\]
These counts are in $A$ gates, not unit-cost $R$ operations.
Moreover
\[
 2bt<12bs\le12(v+b)\le24v,\qquad
 2b<6\sqrt v\le v^{2/3}\quad(v\ge v_0).
\]
Induction using $S(v)\le K v(\log_2v)^8$ above the direct branch
contracts the recursive term by
\[
 24(2/3)^8=6144/6561<1.
\]
The $O(v\log v)$ toll is absorbed by increasing a fixed $K$;
the finite direct branch fixes an effective base constant. This
proves the size bound. The logarithms of child lengths contract
by $2/3$, so the depth recurrence sums to $O(\log(2v))$.

All nonrecursive maps are the printed linear maps; recursive
leaves multiply a left-operand form by a right-operand form.
This proves a formal bilinear identity over $A$, without
$x^2=x$. For generation, compute $b,s,t$ by integer loops, emit
radix-three butterflies with rotation/fold source indices, and
recursively emit the $t$ children. Each descriptor has
$O(\log v)$ address bits; source indices can be obtained by
schoolbook arithmetic and sequential list scans in polynomial
time. This finite recursive printer is uniform, with no root
table, factorization, or coefficient-dependent branch.
\end{proof}

Here is a second explicit way to see that specialization leaves only
linear gates. After fixing the kernel, attach to each gate its
constant value $c_g$ and its homogeneous degree-one part $\ell_g$
in the remaining input indeterminates. An input $x_i$ has pair
$(0,x_i)$, a fixed scalar $a$ has $(a,0)$, and the rules are
\[
 (c,\ell)+(d,m)=(c+d,\ell+m),\qquad
 (c,\ell)(d,m)=(cd,c\,m+d\,\ell).
\]
The constants are prepared once. Only additions and multiplication
by these fixed constants remain in the linear DAG. The rules are
formal truncation modulo total degree at least two; they use no
Boolean relation $x_i^2=x_i$. Since fixed-kernel convolution is
homogeneous linear, its output is exactly the degree-one part.
The transformation changes size and depth by constant factors.

Bounded-degree copy trees, including copies of inputs and outputs,
cost size proportional to the number of use edges. Replacing each
gate's fan-out by its balanced tree adds at most $O(\log v)$ depth
per original circuit layer. The deliberately loose linear-DAG bounds
are therefore
\[
 g(v)=O(v\log^8(2v)),\qquad H(v)=O(\log^2(2v)).
\]
All wire and scalar descriptors have polynomial-time generation in
$v$ and the explicit coefficient dimension. Kernel-specific constant
propagation is charged once in that setup, with the actual coefficient
words retained. This construction is
independent of the faster multiplier being proposed.

\paragraph{Replay on arbitrary dirty coordinates.}
Let a linear DAG implement $M:A^e\to A^a$, and assign a distinct
borrowed coordinate to each of its $g$ gates. The physical source
coordinates $x$, target coordinates $y$, and borrowed coordinates $z$
are pairwise disjoint. Every output is a designated wire; any additions
needed to form it and any input/output copy trees are included in the
counted DAG. Thus the output-read edges also have bounded degree.
In the notation below, $F,D$ name these actual port reads, rather than
uncosted dense matrix applications.
Replace gate assignment
by addition of its predecessor form into that coordinate, in
topological order. For initial borrowed vector $z$ this sweep has
form $z\mapsto Uz+Tx$, where $U$ is unit lower triangular.
If the prescribed reads are $Fz+Dx$, zero initialization would compute
the DAG, so $FT+D=M$.
Sweep, add these reads to $y$, and undo the sweep. The net increment
is $Mx+FUz$, while $z$ is restored. Now sweep with every direct edge
from $x$ omitted, including edges to input-copy-tree roots, and add
only the borrowed-coordinate read $FUz$ to $y$,
and undo again. In characteristic two the two $FUz$ terms cancel.
This realizes $(x,y,z)\mapsto(x,y+Mx,z)$ for every initial $z$,
including coordinates correlated with current sources and targets.

At each DAG depth destinations are fresh gate coordinates and
predecessors are earlier ones. Bounded fan-in and fan-out give a
bounded-degree graph of shear edges. A constant edge coloring makes
each color disjoint pairs. Apply its colors before advancing to the
next depth; undo in reverse depth and color order. The output-port
edges into the distinct $y$ coordinates have the same constant-color
bound. In the second read omit every direct $x$-to-output edge as well.
Zero-coefficient edges may be omitted throughout.
The four sweeps and two output rounds
thus use $O(H+1)$ shear layers, on exactly $e+a+g$ coordinates.
No clean gate values or fresh zero coordinates are assumed.
Other parent coordinates can be selected as borrowed gates, even if
they contain data from an earlier chunk operation; the replay restores
them exactly. A source/target alias is outside this contract.

\paragraph{Explicit rectangular displacement reconstruction.}
\begin{lemma}[Exact-width unit Toeplitz word]
\label{lem:exact-width-toeplitz}
Over any commutative $\F_2$-algebra $A$, a unit lower triangular
Toeplitz matrix $T_v$ has a word on exactly its original $v$ coordinates
of $O(\log^{19}(2v))$ layers. Each layer consists of shears
$(x,y)\mapsto(x,y+\mu x)$ on disjoint ordered pairs.
No clean coordinate is required. Descriptors and coefficients can be
prepared by polynomially many operations on indices, finite lists,
and public coefficient-algebra elements.
\end{lemma}
\begin{proof}
Write $(T_v)_{ij}=h_{i-j}$ for $i\ge j$, with $h_0=1$.
Its reciprocal sequence is prepared by $g_0=1$ and
$g_k=-\sum_{j=1}^k h_jg_{k-j}$; no coefficient inversion is needed.
For $v\ge2$ split at $s=\lfloor v/2\rfloor$ and put $r=v-s$.
The matrix has blocks
\[
 T_v=\begin{pmatrix}T_s&0\\B&T_r\end{pmatrix}.
\]
After applying the two diagonal children in place, a cross update
$E=BT_s^{-1}$ gives this matrix exactly. Its entries are
\[
 E_{ij}=\sum_{\nu=j}^{s-1}h_{i-\nu}g_{\nu-j}.
\]
For consecutive row/column chunks, the interior difference is
\[
 E_{ij}-E_{i-1,j-1}=-h_{i-s}g_{s-j}.
\]
Let $M$ be that rectangular chunk and let $J_a,J_e$ be lower shifts.
Its difference $\mathcal D=M-J_aMJ_e^T$ is the displayed outer
product plus first-row and first-column corrections, hence the sum
of three explicit outer products. More explicitly, let the chunk's
absolute row and column indices start at $i_0\ge s$ and $j_0<s$,
and use zero-based local indices. Set
\[
 v_i=-h_{i_0+i-s},\qquad w_j=g_{s-j_0-j},\qquad
 r_j=M_{0j}-v_0w_j,
\]
and set $c_0=0$, $c_i=M_{i0}-v_iw_0$ for $i>0$. Then, with $e_0$
the first standard vector of the appropriate length,
\[
 \mathcal D=vw^T+e_0r^T+ce_0^T.
\]
The interior follows from the displayed difference identity, the
second term corrects the complete first row, and the third corrects
the remaining first column without counting $(0,0)$ twice. Every
entry is given by finite sums and products of the prepared Toeplitz
and reciprocal coefficients. This uses no field rank test.
Nilpotence gives
$M=\sum_{l\ge0}J_a^l\mathcal D(J_e^T)^l$.
For one term $vw^T$, its action on $x$ is
\[
 b_l=\sum_{j=l}^{e-1}w_{j-l}x_j,\qquad
 c_i=\sum_{l=0}^{e-1}v_{i-l}b_l.
\]
Here $v_k=0$ outside $0\le k<a$.
For literal convolution windows put
$W(Z)=\sum_{j<e}w_jZ^j$, $X^{\rm rev}(Z)=\sum_{j<e}x_jZ^{e-1-j}$
and $V(Z)=\sum_{i<a}v_iZ^i$. The polynomial with coefficients $b_l$ is
\[
 \operatorname{rev}_e\bigl((W X^{\rm rev})\bmod Z^e\bigr),
\]
and the output polynomial is its product with $V$, reduced modulo $Z^a$.
Here $\operatorname{rev}_e P=Z^{e-1}P(Z^{-1})$ at fixed padded width.
The two ordinary products have output lengths at most $2e-1$ and
$a+e-1$, respectively. If $a,e\le b$, both are at most $2b-1$.
Their padding zeros are formal constants; they introduce no extra
runtime source coordinates. Reversal and selection specify DAG ports.
The three displacement terms therefore use six fixed-kernel convolution
circuits and their sum, with total gate bound $O(b\log^8(2b))$ and
depth $O(\log^2(2b))$ including all intermediate and copy coordinates.

At setup enumerate $b=1,\ldots,v$, partition both sides of the cross
block into consecutive chunks of width at most $b$, and print the
actual DAG for every pair, including ragged final chunks and its
input/output copies. Choose the largest $b$ for which
$e+a+g\le v$ for every pair. For a uniform fixed $C$ the count is
$e+a+g\le2b+Cb\log^8(2b)$. Thus
$b_0=\lfloor c v/\log^8(2v)\rfloor$ passes for sufficiently large $v$
and a sufficiently small fixed $c>0$; the chosen size is at least $b_0$.
If none passes, use the direct triangular shear word, updating targets
in descending order from their earlier source coordinates.
This branch occurs only in a fixed finite range.

For each accepted pair its source chunk lies in $[0,s)$ and its target
chunk in $[s,v)$, hence they are disjoint. Choose the $g$ borrowed
coordinates from their complement inside this same $v$-coordinate
parent. Replay restores that complement before the next pair.
The source half is unchanged by every completed cross update, so the
sum of all chunk updates is exactly $E$.
There are
$O(\log^{16}v)$ chunk pairs, each with $O(\log^2v)$ replay layers.
The diagonal children use only their own disjoint intervals.
Run their matching layers in parallel, padding the shorter word
with identity layers. Their union is again a matching. Thus
\[
 d(v)\le\max(d(\lfloor v/2\rfloor),d(\lceil v/2\rceil))
       +O(\log^{18}(2v)),\qquad d(v)=O(\log^{19}(2v)).
\]
All coordinates belong to the original $v$-coordinate parent. Preparing
and counting graphs, reciprocal coefficients, outer-product generators
and edge colors uses only polynomial-size finite lists and polynomially
many public coefficient operations. The search is setup for the public
kernel, never a per-payload-record computation.
\end{proof}

The Newton Toeplitz sequence is $h_j=H_j^{-1}$, so $h_0=1$ as required
by Lemma~\ref{lem:exact-width-toeplitz}. Transposing a layer word reverses
its layer order and each shear's source/target direction; this supplies
the upper triangular word on the same coordinates at the same layer
count. This transposition is on $A$-coordinate indices and retains each
coefficient $\mu$. It is a descriptor operation, not a physical array transpose.
Consequently $F_p=ND^{-1}N^T$ uses two such Toeplitz words and a fixed
number of unit-diagonal stages. Synchronize different prime axes by
identity matching layers. With $p\le R_0$, preparing all their words and
actual $B$-coefficient tables has bit cost polynomial in $R_0,D',L_0$,
as in the coefficient setup below. Each matching layer is then replaced
by the fixed local slot word in the main text; these abstract layer
counts do not assign unit bit cost to growing scalars or tape movement.

\paragraph{The actual constant extension.}
Every irreducible factor of $\Phi_{\alpha'}$ has degree
$r=\operatorname{ord}_{\alpha'}(2)$, so $S$ is a product of copies
of $\F_{2^r}$. A factor tensored with $\F_4$ splits into
$\gcd(r,2)$ fields of degree $r_B=\operatorname{lcm}(r,2)$.
Therefore the same Frobenius exponent in the main proof selects
$\mu=1$ in every field component of $B$, with $r_B\le2D'$.
Computing $r$ by repeated multiplication by $2$ modulo $\alpha'$
needs at most $D'$ iterations and has polynomial small setup cost.
Exponentiation for the selector needs $O(r_B)$ ring products.
The fixed $\F_{16}$ codec inside the binary $H$ supplier need not
embed into $B$: the completed $H$ acts identically on each of the
$2D'$ binary lanes and is thus the required $B$-symbol map.

The central $\F_4$ in $B$ and a fixed field codec used inside a
binary routing/H kernel need not have the same intermediate bit
encoding. The replacement is made only at a completed binary kernel:
its target matrix is $H$, $I$ or $J$ on whole symbols and acts
identically on every binary lane. It therefore implements the required
$B$-linear map in the caller's original representation. Outer
multiplications by elements of $B$ are inserted only after that
boundary and always use the separately charged coefficient-ring
arithmetic. Even a central $\F_4$ action on two concatenated $S$
lanes can be implemented by a constant number of aligned word scans;
no growing basis conversion or free bit transpose is assumed.

\paragraph{Actual growing-word arithmetic and unit setup.}
\begin{lemma}[Fixed monic remainder at any input width]
\label{lem:fixed-monic-remainder}
Fix a monic binary polynomial $f$ of degree $d\ge1$ and a public input
width $r\ge d$. After polynomial-$r$ setup, the remainder modulo $f$
of any binary polynomial of degree $<r$ costs $O(M_{\rm old}(r))$
bit steps on the fixed tape model, including all scans and returns.
\end{lemma}
\begin{proof}
If $r=d$, the input is already reduced. Otherwise set $k=r-d>0$.
The reversal $f^*=X^df(X^{-1})$ has constant coefficient one by
monicity. Prepare $h=(f^*)^{-1}\bmod X^k$ once, for example by its
coefficient recurrence with elementary polynomial-time tape arithmetic.
For the dividend $w$ put
\[
 w^*=X^{r-1}w(X^{-1}),\qquad
 q^*=(w^*h)\bmod X^k,\qquad
 q=X^{k-1}q^*(X^{-1}),\qquad R=w-qf.
\]
All reversals have the displayed fixed padded widths. To prove the
formula, write the unique monic division $w=Qf+R_0$ with
$\deg Q<k$ and $\deg R_0<d$. Reversal gives
$w^*=Q^*f^*+X^{r-1}R_0(X^{-1})$; the second term is divisible by
$X^k$. Multiplying by $h$ modulo $X^k$ therefore recovers $Q^*$,
so $q=Q$ and $\deg R<d$.

The first product may truncate $w^*$ to its first $k$ coefficients.
It has two inputs of length $k$; the other product has input lengths
$k$ and $d+1$, each at most $r$. Thus two calls to the independent
word supplier cost $O(M_{\rm old}(r))$, and all reversals, cropping,
XOR and head returns add $O(r)$. There is no per-dividend Euclidean
algorithm. The prepared inverse length is $k$, even when $d\ll r$.
\end{proof}

Put $d=D'$ and $f=\Phi_{\alpha'}$. Store a $B$ symbol as the two
consecutive $d$-bit $S$ words in $a+b\eta$. If
$P=ac$, $Q=be$, and $T=(a+b)(c+e)$ in $S$, then
\[
 (a+b\eta)(c+e\eta)=(P+Q)+(T+P)\eta.
\]
Three binary polynomial products modulo fixed $f$, supplied by
Lemma~\ref{lem:old-word}, with linear XOR scans, therefore suffice.
Their degree-$\le2d-2$ outputs use
Lemma~\ref{lem:fixed-monic-remainder} at input width $2d-1$.
Its reciprocal is prepared at length $d-1$ for $d\ge2$; $d=1$
is direct. This proves $O(M_{\rm old}(d))$ application time with
no dense $d$-by-$d$ matrix operation.
The runtime word supplier is classical integer SS with carry-free
binary packing; it is separate from the formal characteristic-two
convolution circuit used to construct the local linear DAG above.

The involution $\eta\mapsto\eta+1$ has norm
$N(a+b\eta)=a^2+ab+b^2\in S$. For a unit $x=a+b\eta$, its norm is
an $S$ unit: it is a product of units in $B$, and an inverse in
the free extension $B=S\oplus S\eta$ supplies an inverse in $S$
by taking its constant coefficient. Thus
\[
 x^{-1}=((a+b)+b\eta)\,N(x)^{-1}.
\]
Polynomial extended Euclid with $f$ suffices for the norm inverse,
without factoring $B$. Such inverses are prepared once in polynomial
time in the small dimensions. The selector exponent $2^{r_B}-1$
also needs only $O(r_B)$ binary squaring/product steps, with $r_B\le2d$.
For example,
from $z=\mu+1$ repeat $z\gets z^2(\mu+1)$ exactly $r_B-1$ times.
No possibly nonunit $\mu+1$ is inverted. The one-time polynomial
setup is separate from every runtime scalar application.

\paragraph{Scalar lists on fixed tapes.}
At a synchronized local slot, each axis has a public table
$a_i(x_i)\in B^*$ in that slot's current coordinate order; the
tensor diagonal is $\prod_i a_i(x_i)$. Pairing and coefficient
choices are axis-local. Jointly controlled CRT coefficients use
their different generator and are not covered by this factorization.

Let $M_j=\prod_{i\le j}n_i$, $M_0=1$, with axes in physical
lexicographic order. Start with the one-word list $[1]$.
For every partial product $c$ on the current list, scan the entire
next local table and print
$ca_j(0),\ldots,ca_j(n_j-1)$. Rewind that local table, then advance
to the next $c$. Its rewind costs the size of the child block just
printed. Alternate two expansion tapes and erase the previous list
at the end of each level. This keeps the produced list in the
current native order and uses a fixed tape set.

Level $j$ uses $M_j$ growing-ring products and
$O(M_j(d+L))$ copying, counters and returns. Here $2d\ge L$, and
all $n_i\ge2$, so
\[
 \sum_j M_j<2M,\qquad
 \cost(\hbox{generate and apply one tensor diagonal})
       =O(MM_{\rm old}(d)).
\]
Apply the final list to the payload stream in lockstep and erase it.
Prepare or select each local table once per axis/slot in the
polynomial small setup, using one shared archive and one reusable
active-table tape; there is no tape per axis or indexed archive lookup per
payload record. Rebuild the list from the next slot's tables after
completed relabeling, $H$ and $J$ calls. Transporting an old scalar
list and presuming alignment is not part of the program.

\section{Equal-chunk interchange from the finite rational-frame motif}
\label{app:swap}

The finite rational-frame network and array-interchange strategy are
adapted from OpenAI's Sections 3--4~\cite{swap}. The explicit reconstruction
below supplies the parameters and tape accounting used in this paper.

This appendix constructs, rather than assumes, a fixed-tape interchange
algorithm.  For every binary width $u\ge1$, positive integers $P,G,B$,
and arbitrary $B$-bit record payload, the machine takes an array in
lexicographic order on
\[
 [P]\times[2^u]\times[G]\times[2^u]\times[B]
\]
to the array whose two $u$-bit fields are exchanged, retaining $P,G,B$
and the internal order of each record.  Its time is
\[
                        O(Vu^\tau),
\]
where $V=PGB2^{2u}$ is the entire stored volume and one fixed exponent
$\tau=1-2^{-36}<1$ suffices.  All counters, fixed-table preparation,
copying, row padding, recursive returns, and cleanup are included.  No
unverified global multiplication algorithm or equal-chunk interchange
oracle is a premise.

\subsection{The finite scalar routing network}

Let $\mathcal T=\binom{[g]}3$ with $g=48$.  Two triples are adjacent when
their intersection has size one.  Put
\[
 v=\binom{48}{3}=17296,\quad z=3\binom{45}{2}=2970,
 \quad N=v^3,\quad I=3v^2.
\]
For $S,T\in\mathcal T$, define the binary maps
\[
 \begin{aligned}
 V_{(S,T),T'}&=[T=T']\,[|S\cap T|=1],\\
 J_{S,(S,T)}&=[|S\cap T|=1],\\
 G_{i,T}&=[i\in T],\qquad R_{S,i}=[i\in S].
 \end{aligned}
\]
Then $(JV)_{S,T}$ is one exactly when $|S\cap T|=1$, and
$(RG)_{S,T}=|S\cap T|\pmod2$.  If $S=T$, the intersection size is
three and only $RG$ contributes one.  For distinct triples with
intersection one both terms contribute and cancel in $\mathbb F_2$;
for intersection zero or two both are zero.  Hence
\[
                              RG+JV=I_v.
\]

One invocation has source $x$, target $y$, arbitrary side scratch $a$,
and arbitrary central scratch $c$.  The chronological eight rows are
\[
\begin{aligned}
 y&\gets y-Ja,& y&\gets y-Rc,& a&\gets a+Vx,& c&\gets c+Gx,\\
 y&\gets y+Rc,& y&\gets y+Ja,& c&\gets c-Gx,& a&\gets a-Vx.
\end{aligned}
\]
For initial scratch $(a_0,c_0)$ the target increment is
\[
 -Ja_0-Rc_0+R(c_0+Gx)+J(a_0+Vx)=(RG+JV)x=x,
\]
and the final scratch is exactly $(a_0,c_0)$.  The signs are harmless
over $\mathbb F_2$, but writing them makes the restoration identity
visible.  Thus the invocation adds its source to its target for every
initial scratch value.

Use three stages.  The first is forward from bank $X$ to bank $Y$, the
second is the inverse schedule from $Y$ to $X$, and the third is forward
from $X$ to $Y$.  On arbitrary bank contents,
\[
 (x,y)\longmapsto(x,y+x)\longmapsto(-y,y+x)\longmapsto(-y,x).
\]
Over $\mathbb F_2$ this exchanges $X,Y$.  Each triple coordinate
participates in each stage once; invocations have disjoint scratch
roles.  All scratch roles are restored.  The resulting finite role
permutation $\rho$ exchanges the $N$ data roles and fixes every scratch
role.  It is the scalar-routing network used below.  The finite incidence
identity and scratch schedule are also recorded in
\path{supporting/SCALAR_MOTIF_SYMBOLIC_AUDIT.md} and
\texttt{finite\_phase\_appendix.tex}.

\subsection{Rational frames and exact child count}

Address arithmetic is over $\mathbb Q$ first and then over
$\mathbb Z/(q^b)$ for a fixed odd prime $q$.  It is distinct from the
payload network over $\mathbb F_2$.  On $F=\mathbb Q^{48}$ use the
nondegenerate symmetric form with Gram matrix
\[
                         B_0=I-J/9.
\]
The all-one line has eigenvalue $1-48/9=-13/3$, and its orthogonal
complement has eigenvalue one, so the form is nondegenerate.  For
$U\in\mathcal T$, let $t_U=\mathbf1_U$.  Then
\[
 \langle t_U,t_V\rangle=|U\cap V|-1,\qquad
 \langle t_U,t_U\rangle=2.
\]
Thus $P_U$, the orthogonal projection onto $\langle t_U\rangle$, is a
well-defined rational rank-one projection.  Adjacent triples give
orthogonal lines.  The common-label table for the finite motif follows
from these orthogonal direct-sum decompositions; a derivation is in
  \path{supporting/ntt_base_five/finite_motif_rank_closed_form_audit.md}.

At one frame stage, let $A$ be the preceding tensor label, $P$ its
prefix line, $B=P^\perp$, and $L_U=\langle t_U\rangle$.  The next center
label satisfies
\[
 (P\otimes L_U)^\perp
   =(B\otimes F)\oplus^\perp(P\otimes L_U^\perp).
\]
Side inclusions use $L_U\perp L_V$ for adjacent $U,V$.  In the fixed
physical gate order the common subspace labels are
\[
\begin{array}{c|l}
0&B\otimes L_U\\
1&B\otimes F\\
2&(B\otimes F)\oplus^\perp(P\otimes L_U)\\
3&A\otimes F\\
4&B\otimes F\\
5&(B\otimes F)\oplus^\perp(P\otimes L_U^\perp)\\
6&A\otimes F\\
7&A\otimes F.
\end{array}
\]
Every label is tensored with the fixed suffix tensor line of the
invocation. A side wire named by its physical Y/X neighbor pair $(S,T)$
meets rows $0(S),2(T),5(S),7(T)$; its middle inclusion is
$L_T\subset L_S^\perp$. A center wire meets rows $1,3,4,6$.
Data X meets $2,3,6,7$ and data Y meets $0,1,4,5$.
The second stage uses the reversed schedule with logical banks exchanged,
so its ordered side-pair names are reversed to retain this physical
convention. Scratch input frames are zero, scratch output frames are
$I_m$, and the final data X/Y frames are $I_m$ and $I_m-P_U$.
These finite inclusions and wire incidences use no phase-kernel,
weighted-mask, H2, or interchange algorithm as a premise.

In the fixed
eight-role order the label dimensions are
\[
 a-1,\quad 48(a-1),\quad48(a-1)+1,\quad48a,\quad
 48(a-1),\quad48a-1,\quad48a,\quad48a,
\]
where $a=\dim A$.  The only decreases are the central transitions
$48a\to48(a-1)$, each losing 48 dimensions.  With the fixed counts
\[
 W=2N+I(vz+48),\qquad m=48^3=110592,
 \qquad L=I\cdot48^2,
\]
the signed endpoint dimension telescope and twice the downward variation
give the sum of comparable-edge ranks $Wm-2N+2L$.

\paragraph{Complete physical edge ledger.}
The inverse middle stage runs logical rows $7,6,\ldots,0$ with source
Y and target X. Its physical gate incidences are nevertheless
Y-side, Y-center, X-side, X-center, Y-center, Y-side, X-center, X-side,
the same chronological order as the other stages. Thus the common-label
table is used in its printed order, with the side-pair names reversed;
one must not reverse the table along with the logical row numbers.
Put $d=48$ and $a=d^{j-1}$ at stage $j$. Suppressing the suffix line,
the complete scratch dimension paths, including their endpoints, are
\[
\begin{array}{c|l|c}
 \hbox{physical wire}&\hbox{dimension path}&\hbox{downward loss}\\ \hline
 \hbox{side}&0,\ a-1,\ d(a-1)+1,\ da-1,\ da,\ m&0\\
 \hbox{center}&0,\ d(a-1),\ da,\ d(a-1),\ da,\ m&d .
\end{array}
\]
These are geometric inclusions: the middle side edge uses
$L_T\subset L_S^\perp$, whose nondegenerate residual has dimension
$d-2$; the center's downward residual is $P\otimes F$.
At an interstage boundary X enters as $A\otimes L_T$, contained in
$(B\otimes F)\oplus(P\otimes L_T)$, while Y enters as
$B\otimes L_S$, exactly its next first label. The physical X path is
therefore nondecreasing from a formal positive input line to dimension
$m$, and Y is nondecreasing from zero to dimension $m-1$.
Their ordinary edge-rank sums are both $m-1$.
Each side contributes $m$ and each center $m+2d$.
Since there are $I\cdot48$ centers, summing gives directly
\[
 2N(m-1)+(W-2N)m+2I\cdot48d=Wm-2N+2L.
\]
The actual negative X source, treated next, replaces its formal
zero-rank first edge by one rank-one edge; it does not alter the
ordinary interstage inclusions.

At each of the $N$ exceptional $X$ sources choose input frame $-P_U$;
the first gate frame is $+P_U$.  Their difference is
$2P_U$, of rational rank one.  Since the address modulus will have odd
characteristic, this is not zero.  These $N$ source corrections add $N$
recursive children, so the full child count is
\[
                         s=Wm-N+2L.
\]
For every routed role $w$, the endpoint labels obey
\[
                  M_{\rm out}(\rho(w))-M_{\rm in}(w)=I_m.
\]
For an $X$ data role the endpoint difference is
$P_{U^\perp}-(-P_U)=I_m$; the other data and scratch endpoint pairs
are the corresponding zero/full or complementary-label pairs in the
finite common-label table.  For comparable labels, the rank of the
projection difference is the dimension change.  The only equal-dimension
exception is the $N$ first $X$ edges, already counted as $2P_U$ of rank
one.  Thus no source correction is hidden by the endpoint identity.
At each role edge with label difference $A$, implement the corresponding
address-frame change.  The scalar routing sends role $w$ to $\rho(w)$;
the displayed endpoint identity then leaves the same target shear on
data and scratch roles alike.

The exact fixed counts reduce to
\[
 \frac WN=\frac{9633881}{1081},\qquad
 \frac LN=\frac{432}{1081},\qquad
 1-\frac{s}{Wm}=\frac{217}{9633881\cdot110592}.
\]
The integer inequality
\[
 217\cdot2^{36}>12\cdot9633881\cdot110592
\]
and $m<2^{17}$ give $\ln m<12$ and
$1-s/(Wm)>12\cdot2^{-36}$.  From $e^{-x}>1-x$ it follows that
\[
                      s/W<m^\tau,\qquad \tau=1-2^{-36}.
\]
All constants here are fixed.  These are exact rational counts, not a
finite-size matrix-rank inference.

\subsection{Choosing one odd address prime}

There are finitely many rational label matrices, projections, and frame
differences.  For each edge matrix the construction below gives a fixed
factorization into lower-triangular matrices and a partial permutation.
Collect the denominators of all entries and the numerators and
denominators of every nonzero diagonal unit in these factorizations and
their inverses.  This is a finite set of nonzero integers.  Choose one
odd prime $q$ dividing none of them.  Such a prime exists by Euclid's
argument; trial division through primes finds one after finitely many
steps.  This is fixed finite setup, independent of $b,u,P,G,B$.  The
prime and the rational matrices may equivalently be constants of the
finite machine.

All rational matrices now reduce to $\mathbb Z/(q^b)$ for every $b$.
The chosen triangular diagonal entries remain units.  Each fixed
factorization remains an exact identity, with the same $a$ partial
permutation pivots.  No data-dependent division or prime oracle is used.
Because $q$ is odd, the exceptional source matrix $2P_U$ still has its
one nonzero pivot and costs one child swap.

\subsection{Fixed-radix tape operations}

The following operations use fixed $q$ and fixed tape families.

\paragraph{Counters and records.}
Use little-endian ripple counters with the head returning to the low
end. Over a run of $A$ increments, or a run of $A$ decrements, of an $r$-bit
counter from an arbitrary initial state, bit $j$ is visited
$O(1+A/2^j)$ times, so the total cost is $O(A+r)$.
The width term is necessary: one increment from $2^r-1$ can visit
all $r$ bits. For a nonempty range of $n$ values, use its own
$r=\max\{1,\lceil\log_2(n+1)\rceil\}$ active counter width and
initialize/reset it at the prescribed endpoint. A full scan,
initialization and return then cost $O(n+r)=O(n)$.
To test termination without a full-width comparison after each step,
initialize a countdown to $n-1$, execute the loop body, then
decrement and stop precisely on borrow out of its marked $r$-bit
interval. The last decrement from zero performs the end test
while visiting those bits. This gives exactly $n$ bodies, including
$n=1$. A needed forward index uses a separate amortized counter.
Range bounds are prepared once in their own canonical widths;
each reset copies only that bound, not a full ambient-width register.
Stale cells outside this active interval are not rescanned at every
reset. With a fixed number of nested ranges, charge each inner reset
to its ensuing nonempty scan. Copying a $B$-bit record thus costs
$O(B)$ including its counter initialization and returns, even for
$B=1$. Recursion saves counters/descriptors at the current stack top;
it does not cross an ancestor to fetch a length. Ordinary schoolbook
shape arithmetic has a fixed polynomial cost in descriptor bitlength;
the retained-range bound below pays it at every recursive call.

More explicitly, traversing a fixed-depth product of nonempty ranges
$n_1,\ldots,n_d$, including a final bit-record range if necessary,
has counter cost bounded by
\[
 O\!\left(\sum_{i=1}^d
       \Bigl(\prod_{j<i}n_j\Bigr)
       (n_i+\log(2n_i))\right)
   =O\!\left(d\prod_{i=1}^d n_i\right).
\]
For the current fixed matrix dimension, $d$ is fixed, even when
prefix, gap and suffix ranges are arbitrarily large. All array
bits remain raw payload: counters and boundary markers are on
work tapes, with no growing header added to each record. Rewinding
the fixed number of current split tracks costs their combined
volume and stops at their current stack boundaries.

\paragraph{Digit moves and pointwise maps.}
Split a stream into $q$ tracks by one selected radix-$q$ digit, preserving
order on each track, and merge in the desired digit order.  A single
digit move costs $O(V)$.  For a fixed permutation of two digit positions,
split into $q^2$ tracks and merge according to the two preimage digits;
this also costs $O(V)$.  A fixed pointwise finite-alphabet map on aligned
tracks costs $O(V)$.  The number $q^2$ is large but fixed, not input
dependent.

\paragraph{Rotations controlled by earlier fields.}
Let the target field have $Q=q^b$ values and follow every field used to
compute its shift.  Within each fixed earlier-address fiber, it consists
of $Q$ consecutive blocks of suffix records.  For shift $a\in[Q]$,
copy the first $Q-a$ blocks to one stream and the last $a$ to another,
then write the latter before the former.  This realizes
$y\mapsto y+a\pmod Q$ in $O(V)$ record motion.  Computing the offset by
schoolbook arithmetic on the fixed number of $b$-digit controls costs
$\operatorname{poly}(b)$ per fiber.  Each fiber contains $q^b$ blocks,
so this metadata is absorbed into its own volume.  Arbitrary spectators
and payload remain inside each block.
The enumeration retains separate counters for the selected controls.
Large spectator counters are advanced or reset by the preceding
product-scan rule; their full values are not decoded or copied into
the offset calculation. Thus the bound is $O(V)$ even for $b=1$,
with arbitrarily large prefix, intervening gap, or suffix.
This argument relies on reading only the fixed number of selected
$b$-digit fields. It does not authorize an offset depending on the
entire ambient address at the same cost.

\paragraph{Fixed rational units.}
For a fixed positive integer $c$ coprime to $q$, partition $y\in[Q]$
into its $c$ monotone pieces indexed by $j=\lfloor cy/Q\rfloor$.
On a piece the output is $y'=cy-jQ$ in increasing order.  Since
$Q$ is invertible modulo $c$, for each output value $y'$ the unique piece
is determined by $j\equiv-y'Q^{-1}\pmod c$.  Split by these fixed $c$
classes and merge their increasing streams in output order.  This is
$y\mapsto cy\pmod Q$ in $O(V)$, including piece endpoints and counters.
Prepare the endpoints and $Q^{-1}\bmod c$ in polynomial-$b$ time;
the output selector needs only $y'\bmod c$, a fixed finite-state
counter updated once per suffix block. It performs no full $b$-bit
multiplication or global-address decoding at each record.
The same residue split on a supplied output stream constructs the inverse
permutation directly.  Negation modulo $Q$ fixes zero and reverses the
other $Q-1$ blocks; a stack reversal plus reversal within each block costs
$O(V)$.  A fixed rational unit $c/d$ is the composition of integer
multiplication by $|c|$, inverse multiplication by $|d|$, and a possible
negation.  The prime $q$ avoids $c,d$.

\subsection{Lower-triangular frames and rank-many child swaps}

For a fixed rational matrix $A$ of rank $a$, eliminate its topmost
nonzero row using the rightmost nonzero pivot.  Clear entries to the
left by lower-triangular column operations and entries below by
lower-triangular row operations.  Remove the isolated pivot row and
column and continue in the remaining ordered indices.  This yields
\[
                         A=E_1\Pi E_2,
\]
where $E_1,E_2$ are nonsingular lower-triangular matrices and $\Pi$
has exactly $a$ ones in distinct rows and columns.  Reversing the
elimination proves the factorization; nonsingularity of $A$ is not
required.
Concretely, if the pivot is in row $i_0$, column $j_0$, clearing a
column $j<j_0$ adds a multiple of column $j_0$ to column $j$,
a lower-triangular right operation. Clearing a row $i>i_0$ adds a
multiple of row $i_0$ to row $i$, a lower-triangular left operation.
Normalize the pivot by a nonzero diagonal scale. Previously isolated
pivot rows and columns are unchanged, since their entries in the
remaining submatrix are zero. Retaining the original order on the
remaining indices preserves triangularity. Thus $LAR=\Pi$ with
invertible lower-triangular $L,R$, and $E_1=L^{-1},E_2=R^{-1}$.
The partial permutation need not be diagonal.

\begin{lemma}[Signed factorization over every address power]
\label{lem:swap-signed-module}
For each actual directed rational frame difference $A=M_j-M_i$ of
rank $a$, the chosen fixed prime gives, for every $b\ge1$,
\[
 A_b=E_{1,b}\Pi E_{2,b}\quad\hbox{over }R_b=\mathbb Z/(q^b),
\]
with both triangular factors invertible. The image and kernel of $A_b$
are free direct summands of ranks $a$ and $m-a$. The same $a$ unit
pivots implement the signed address shear; no field assumption on
$R_b$ or invertibility of a data coordinate is required.
\end{lemma}
\begin{proof}
The finite rational identities reduce entrywise since $q$ avoids every
entry denominator. Each triangular diagonal remains a unit by the
same fixed exclusion, so both factors are invertible over every $R_b$.
Multiplication by those invertible factors identifies the image and
kernel with those of $\Pi$, whose nonzero entries are literal ones.
In particular these are coordinate direct summands of the stated ranks.
This is an explicit factorization certificate over a ring, rather than
a rank test over a field.

Factor the \emph{signed} matrix $A$ before reduction. Its negative
can equally be factored as $(-E_1)\Pi E_2$, absorbing the sign in
unit diagonals. For the source correction $2P_U$, the factor two
is likewise absorbed in one triangular factor; $q$ is odd.
It retains one pivot and one child. No sign or factor two is reduced
using the separate characteristic-two payload arithmetic.
\end{proof}

A lower-triangular coordinate change within a bank is performed in
descending coordinate order.  At target coordinate $i$, the controls
with index $j<i$ still have their original values and physically precede
the target.  Apply its fixed diagonal unit, then add the fixed linear
combination of those earlier coordinates by later-target rotations.
This costs $O(V)$ per coordinate; the dimension $m$ is fixed.

Define, for vectors of $m$ radix-$q^b$ digits,
\[
 \Phi_A(H,D)=(H+AD,D).
\]
Its induced forward record pushforward is
\[
 (\mathcal P_A f)(H,D)=f(H-AD,D),\qquad
 \mathcal P_A\mathcal P_B=\mathcal P_{A+B},\qquad
 \mathcal P_A^{-1}=\mathcal P_{-A}.
\]
Here the coordinate signs are in $R_b$; pointwise payload addition
remains XOR over $\mathbb F_2$.
To implement it, first change $D$ by $E_2$, change $H$ by $E_1^{-1}$,
add $\Pi D$ to $H$, undo the $H$ change by $E_1$, and finally undo the
$D$ change by $E_2^{-1}$.  The result is
$H+E_1\Pi E_2D=H+AD$ with the original $D$ restored.  The triangular
operations are linear volume.
For each one in $\Pi$, the cross-bank shear $H_i\gets H_i+D_j$ is
implemented in the physical order (earlier field $H_i$, later field
$D_j$) by
\[
       D_j\gets D_j+H_i;\qquad
       \operatorname{swap}(H_i,D_j);\qquad
       D_j\gets H_i-D_j.
\]
On $(h,d)$ this is $(h,d)\mapsto(h+d,d)$.  The first and last
operations target later fields and cost $O(V)$.  The middle operation is
one recursive interchange of two $b$-digit fields across arbitrary gap
and suffix.  Thus $\Phi_A$ costs exactly $a$ recursive child swaps plus
$O(V)$ other work.  Its inverse is $\Phi_{-A}$ and has the same bound.
Lemma~\ref{lem:swap-signed-module} justifies these same fixed factors
for singular, negative and source-correction matrices at every $b$.

\subsection{Exact power widths, role frames, and tape recurrence}

First let each chunk have width $e=m^k$ radix-$q$ digits.  Keep an
independent row field of size $R$ divisible by $W^k$, before both chunks;
arbitrary prefix, gap, suffix, and spectators between the row and first
chunk are retained.  For $k=0$ the two $q$-digit chunks are exchanged by
the fixed split/merge construction in $O(V)$.

For $k>0$, split each chunk into $m$ consecutive width-$b$ fields,
$b=e/m$, and denote their vectors by $H,D$.  Split whole rows cyclically
into $W$ role streams inside each earlier prefix, stream $w$ taking rows
$r=Wa+w$.  Each role has the full $H,D$ ranges and all spectators, and
row count $R/W$ divisible by $W^{k-1}$.  On each finite scalar-network
edge from label $M_i$ to $M_j$, apply $\Phi_{M_j-M_i}$ to its role stream,
then perform the specified pointwise XOR gate on aligned payloads.

For input array $f_w$ on role $w$, use the logical interpretation
$\mathcal P_{-M_{\rm in}(w)}f_w$.  This is a proof invariant, not a
preprocessing pass.  The frame change on an edge is common to its aligned
XOR operands and commutes with pointwise addition.  The scalar network
routes role $w$ to $\rho(w)$; its physical output there is
\[
\mathcal P_{M_{\rm out}(\rho(w))-M_{\rm in}(w)}f_w
       =\mathcal P_I f_w.
\]
Equivalently, with block-diagonal input/output pushforward frames
$\mathcal D_{\rm in},\mathcal D_{\rm out}$ and scalar role
permutation $P$, the entire framed word is
\[
 \mathcal D_{\rm out}P\mathcal D_{\rm in}^{-1}
        =P(I_W\otimes\mathcal P_I).
\]
The endpoint difference is $I_m$ on every routed role, including dirty
side and center roles. A common reindexing commutes with their XOR
gates; no payload value is assumed zero.
Merge output role $\rho(w)$ back into input row class $w$.  This undoes
the role exchange and applies $H\gets H+D$ to every row, including
arbitrary scratch rows, while preserving all spectators.

The finite role network supplies only this earlier-target shear.  Complete
the chunk interchange with the outer later-target operations
\[
 D\gets D-H,\qquad H\gets H+D,\qquad D\gets H-D.
\]
On input $(h,d)$ the result is $(d,h)$.  Only the middle operation uses
the framed role network.  The first and last operations are later-target
rotations and fixed unit maps, each costing $O(V)$.  Thus the number of
recursive children remains exactly $s$; it is not multiplied by three.

The finite edge matrices have ranks summing to $s$. A pivot child is a
direct call to the exact power-width entry. More explicitly, its parent
role shape is
\[
\begin{aligned}
 &[P][R/W][A][H_1]\cdots[H_m]\\
 &[G][D_1]\cdots[D_m][B].
\end{aligned}
\]
Every formal field has width $b=e/m$, with a separate pretarget
spectator range $A$. A pivot using $H_i,D_j$ has child parameters
\[
 A_c=Aq^{b(i-1)},\qquad
 G_c=q^{b(m-i)}Gq^{b(j-1)},\qquad
 B_c=q^{b(m-j)}B.
\]
The call is
$\operatorname{SwapPower}(b,q,P,R/W,A_c,G_c,B_c)$.
The row range remains explicit and is divisible by $W^{k-1}$;
all unselected fields remain in the contiguous spectator products.
Consequently
\[
 P(R/W)A_cG_cB_cq^{2b}
   =P(R/W)AGBq^{2e}=V/W.
\]
This power entry introduces no numeric or row pad. It does not call
the arbitrary-width or binary wrapper. The pivot factorization
therefore gives exactly $s$ recursive calls of volume $V/W$ and
width $e/m$. The
finite number of role gates, fixed-rational operations, XOR scans, row
splits/merges, and setup work contributes $C_1V$ at the parent.  Use
separate fixed tape families and a parking stack: before one child, park
only the other current role streams at the stack top; pass the active
stream to the child; after return, restore siblings by popping and writing
right-to-left.  Clear only current visited intervals.  A descendant never
traverses a parked ancestor.  Since $s,W$ are fixed constants, all parent
parking costs $C_sV$ and does not multiply child time.

With normalized time $F_k$, the exact recurrence is
\[
 F_0\le C_0,\qquad
 F_k\le(s/W)F_{k-1}+C_1.
\]
By the finite count $s/W<m^\tau$, this gives
\[
                         F_k=O(m^{k\tau})=O(e^\tau).
\]
The bound also holds if $s/W\le1$, since the resulting linear sum in
$k$ is at most $O(e^\tau)$.  Parent parking uses $O(V)$ time and
storage at the current level.  At depth $j$, the child volume is
$V_j=V_0/W^j$; the two original full chunk ranges remain present as
targets or spectators, so $V_j\ge q^{2e_0}$ for root width $e_0$.
Since $j\le k=O(\log e_0)$,
\[
 \log V_0=\log V_j+j\log W=O(\log(2V_j)).
\]
Any fixed polynomial of the root descriptor bitlength therefore fits
within $O(V_j)$ at the child.  Its descriptor lists only the current
fields and collapsed spectator intervals, not the chain of ancestor
subdivisions.

\subsection{Arbitrary $q$-radix widths without rounding the core}

Let the two original chunks have width $e$ radix-$q$ digits.  Put
$k=\lceil\log_m e\rceil$.  Choose the least $r_0$ such that
$q^{2r_0}\ge W^k$.  For sufficiently large $e$, $1\le r_0<e$ and
$r_0=O(\log(2e))$; the bounded smaller range uses direct finite
single-digit interchanges and changes only the uniform constant.
Write the first chunk as $h_+h_-$ and the second as $d_+d_-$,
with the four lengths specified separately:
\[
 |h_+|=|d_+|=r_0,\qquad |h_-|=|d_-|=e-r_0.
\]
Thus the plus sign always means high digits and the minus sign
always means low digits. Starting from
\[
 [h_+][h_-][G][d_+][d_-],
\]
first exchange corresponding high digits individually, obtaining
\[
 [d_+][h_-][G][h_+][d_-].
\]
Then move only the $r_0$ digits of $h_+$ left across $h_-$ and $G$
using fixed-radix digit moves, obtaining
\[
 [d_+][h_+][h_-][G][d_-].
\]
This costs $O(Vr_0)$.  The joined row field $(d_+,h_+)$ has size
$R=q^{2r_0}\ge W^k$.  Pad this one independent numeric row range to
\[
 R'=W^k\left\lceil\frac{R}{W^k}\right\rceil,
                   \qquad R\le R'<2R.
\]
It is divisible by every $W^j$ needed in the recursion.  This private
factor below two multiplies the current volume by one constant factor for
the whole arbitrary-width call.  All children receive that same padded
row range and have volume $V'/W$ at the first level (then divide by $W$
per level); the pad is not independently reintroduced at descendant
levels.  Pad rows are initialized to zero.  Completed interchange calls
preserve the row coordinate, so these rows remain zero at the crop
boundary.  Crop the row factor once, after the complete arbitrary-width
core, to recover the original volume.

Write $e-r_0$ in base $m$.  At scale $m^j$ there are at most $m-1$
matching consecutive pieces in each lower chunk.  Interchange each
corresponding pair by the exact power-width routine, with all other
pieces as spectators.  The sum of their costs is
\[
 O(V)\sum_j a_jm^{j\tau}=O(Ve^\tau),\qquad 0\le a_j<m.
\]
No target piece is rounded up to a larger power.  After all pieces,
the field order is
\[
 [d_+][h_+][d_-][G][h_-].
\]
Crop the row padding back to $R$ genuine rows before undoing the high
digit move.  Move $h_+$ after $G$ by the reverse-direction sequence of
$r_0$ digit moves, obtaining
\[
 [d_+][d_-][G][h_+][h_-].
\]
The two original chunks have been interchanged.  The high moves cost
$O(V\log(2e))=O(Ve^\tau)$, and row padding contributes only a fixed
factor below two.  All spectators retain their order.

For completeness, fixed-radix digit moves use $q$-way stream splitting
and merging; their cost is $O(V)$ each and uses a fixed number of tapes.
The $O(\log e)$ piece schedule and shape arithmetic are polynomial in
$\log(2V)$, hence $O(V)$.  The full chunk ranges remain in every
descendant, which is the volume lower bound used for descriptor costs.

Counters and record movement use ordinary sequential tapes.  A
little-endian binary counter changes bit $j$ only $O(1+T/2^j)$ times
over $T$ consecutive increments, so a run of $T$ increments, its reset,
and return cost $O(T+w)$ for width $w$. Use a range-specific active
width as above; do not restart a full ambient-address counter for
each constant-size inner scan. Nested counters and child
descriptors are pushed at the current stack top and only that child's
interval is popped on return.  A $B$-bit suffix record is copied with an
amortized down-counter in $O(B)$, including reset.  At each call there
are only a fixed number of shape parameters; polynomial work in
$\log(2V)$ is $O(V)$ for all positive input volumes after increasing the
uniform constant.  No counter repeatedly scans an ancestor descriptor.

\subsection{Binary fields and final theorem}

For binary width $u$, set $N=2^u$ and choose the least $e$ with
$Q=q^e\ge N$.  Then $Q<qN$.  Pad each of the two numeric address ranges
from $N$ to $Q$, copying the valid rectangle and filling its complement
with zero records.  This nested scan costs $O(V')$ where
\[
 V'=PGBQ^2<q^2V.
\]
Radix-$q$ lexicographic order is numeric order, so no per-record base
conversion is used.  Apply the exact $q$-digit interchange above.  Its
completed coordinate transposition preserves the nested valid rectangle
$[N]\times[N]$; hence invalid padding remains zero and can be cropped
after the full operation.  No crop is performed inside a recursive piece.
Because $e=O(u)$ for fixed $q$, the final time is
\[
                         O(V'u^\tau)=O(Vu^\tau).
\]

All operation families have fixed finite tape sets: the $q^2$ split
tracks, the $W$ role streams, reusable I/O/work tapes, and stack/descriptor
tapes depend only on fixed constants.  Finite rational tables and the
avoiding prime are fixed machine data or are generated by terminating
constant setup.  Counters use ordinary amortized binary increments;
record copies and returns cost their full bit volume.  Thus this is one
uniform fixed-finite-tape construction for all $P,G,B,u$, including
$B=1$.  It proves the local interchange statement without assuming the
source preprint's global integer multiplication algorithm.  It makes no
novelty claim for the finite motif or its routing method.





\begin{thebibliography}{99}
\bibitem{cyclomult}
David Harvey and Joris van der Hoeven.
\newblock Faster polynomial multiplication over finite fields using cyclotomic coefficient rings.
\newblock \emph{Journal of Complexity}, 54:Article 101404, 2019.
\newblock \url{https://doi.org/10.1016/j.jco.2019.03.004}.
\newblock Author manuscript: \url{https://www.texmacs.org/joris/cyclomult/cyclomult.pdf}.

\bibitem{ffnlogn}
David Harvey and Joris van der Hoeven.
\newblock Polynomial multiplication over finite fields in time $O(n\log n)$.
\newblock \emph{Journal of the ACM}, 69(2), Article 12, 1--40, 2022.
\newblock Theorem 1.2 assumes the stated Linnik hypothesis.
\newblock \url{https://doi.org/10.1145/3505584}.
\newblock Author manuscript: \url{https://www.texmacs.org/joris/ffnlogn/ffnlogn.pdf}.

\bibitem{cantor-kaltofen}
David G. Cantor and Erich Kaltofen.
\newblock On fast multiplication of polynomials over arbitrary algebras.
\newblock \emph{Acta Informatica}, 28(7):693--701, 1991.
\newblock \url{https://doi.org/10.1007/BF01178683}.
\newblock Author copy: \url{https://kaltofen.math.ncsu.edu/bibliography/91/CaKa91.pdf}.

\bibitem{schonhage-strassen}
Arnold Sch\"onhage and Volker Strassen.
\newblock Schnelle Multiplikation gro{\ss}er Zahlen.
\newblock \emph{Computing}, 7:281--292, 1971.
\newblock \url{https://doi.org/10.1007/BF02242355}.

\bibitem{swap}
OpenAI.
\newblock \emph{Integer multiplication below $n\log n$}.
\newblock Research preprint, 23 September 2026, Sections 3--4;
Lemma 12 supplies the arbitrary-width equal-chunk interchange statement.
\newblock Citation revision \texttt{4eaee87548250b2682affda7d05d947c05554bae}, inspected 31 October 2026.
\newblock \url{https://github.com/openai/math/blob/4eaee87548250b2682affda7d05d947c05554bae/preprints/Integer-multiplication-below-n-log-n-September-23-2026/build/sections/04-swap.tex}.

\bibitem{local}
OpenAI.
\newblock \emph{An explicit power saving for the exact discrete Fourier transform}.
\newblock Research preprint, 25 September 2026, local Fourier compiler section.
\newblock Citation revision \texttt{4eaee87548250b2682affda7d05d947c05554bae}, inspected 31 October 2026.
\newblock \url{https://github.com/openai/math/blob/4eaee87548250b2682affda7d05d947c05554bae/preprints/An-explicit-power-saving-for-the-exact-discrete-Fourier-transform-September-25-2026/build/sections/local.tex}.

\bibitem{benes}
Meng Zhang and Yi Zhang.
\newblock Arbitrary-size permutation networks using arbitrary-radix switches.
\newblock \emph{Theoretical Computer Science}, 962:113939, 2023.
\newblock \url{https://doi.org/10.1016/j.tcs.2023.113939}.

\bibitem{integer-followups}
Douglas Colkitt and the contributors to \texttt{CrocSwap/integer-mult-bounds}.
\newblock \emph{A sharper exponent for integer multiplication}.
\newblock Research repository, revision \texttt{2485805a6d47f36ea83ddde6327d453ac82d44dc},
accessed 31 October 2026. Contributor-level source notices are retained
in the bundled native derivation. Its selected integer bound retains
its stated framework assumptions.
\newblock \url{https://github.com/CrocSwap/integer-mult-bounds/tree/2485805a6d47f36ea83ddde6327d453ac82d44dc}.

\bibitem{integer-lean}
\texttt{platypii/integer-mult-bounds-lean} contributors.
\newblock \emph{Integer multiplication in Lean}.
\newblock Research repository; the stage-two source used by the companion
odd-field transfer is pinned to \texttt{fefd699f9389af961046997796261333539e5909}.
The claimed global integer multiplier is not a premise here.
\newblock \url{https://github.com/platypii/integer-mult-bounds-lean/tree/fefd699f9389af961046997796261333539e5909}.

\bibitem{HVHL17}
David Harvey, Joris van der Hoeven, and Gr\'egoire Lecerf.
\newblock Faster polynomial multiplication over finite fields.
\newblock \emph{Journal of the ACM}, 63(6), Article 52, 2017.
\newblock \url{https://doi.org/10.1145/3005344}.
\newblock Author manuscript: \url{https://www.texmacs.org/joris/ffmul/ffmul.pdf}.

\bibitem{Harvey26Context}
David Harvey.
\newblock \emph{On the recently proposed $o(n\log n)$ integer multiplication algorithm}.
\newblock Unpublished discussion note, 20 October 2026, Section 4.
\newblock \url{https://web.maths.unsw.edu.au/~davidharvey/research/subnlogn.pdf}.

\bibitem{Frobenius17}
Joris van der Hoeven and Robin Larrieu.
\newblock The Frobenius FFT.
\newblock In \emph{Proceedings of ISSAC}, 437--444, 2017.
\newblock \url{https://doi.org/10.1145/3087604.3087633}.
\newblock Author manuscript: \url{https://www.texmacs.org/joris/ffft/ffft.pdf}.

\bibitem{CodexTool26}
Nicolas Bonamy and Derrick Choi.
\newblock Codex as a platform: build on the open agent harness.
\newblock OpenAI Developers, 19 August 2026. Official documentation
reference for the Codex software used in release preparation and source
inspection; accessed 31 October 2026.
\newblock \url{https://developers.openai.com/blog/codex-as-a-platform}.

\end{thebibliography}
\end{document}



